% Template article for preprint document class `elsart'
% Jin Chen 2013/01/07 
% TESTSUITE/symmetry/hopaper.m, out_pppl.m

a(1,1)= 1
      a(2,2)=1    
            a(3,3)=1    
                  a(4,4)=2    
                        a(5,5)=1    
                              a(6,6)=2    

a(7,1)=1    a(7,2)=1                a(7,4)=1                             a(7,7)=1                                         a(7,11)=1                                                    a(7,16)=1                
            a(8,2)=1                a(8,4)=2                             a(8,7)=3                                         a(8,11)=4                                                    a(8,16)=5                
                        a(9,3)=1                a(9,5)=1                             a(9,8)=1                                          a(9,12)=1                                                                
                                    a(10,4)=2                            a(10,7)=6                                        a(10,11)=12                                                  a(10,16)=20              
                                                a(11,5)=1                            a(11,8)=2                                         a(11,12)=3                                                               
                                                            a(12,6)=2                            a(12,9)=2                                         a(12,13)=2                          a(12,16)=-10 a(12,17)=-2 

a(13,1)=1               a(13,3)=1                           a(13,6)=1                                         a(13,10)=1                                                   a(13,15)=1                           a(13,18)=1
            a(14,2)=1                           a(14,5)=1                                        a(14,9)=1                                                     a(14,14)=1                                       
                        a(15,3)=1                           a(15,6)=2                                         a(15,10)=3                                                   a(15,15)=4                           a(15,18)=5
                                    a(16,4)=2                                        a(16,8)=2                                                     a(16,13)=2                                       a(16,17)=2  a(16,18)=-10
                                                a(17,5)=1                                        a(17,9)=2                                                     a(17,14)=3                                       
                                                            a(18,6)=2                                         a(18,19)=6                                                   a(18,15)=12                          a(18,18)=20

% publish single column
%\documentclass[final, oneeqnum, onefignum, onetabnum]{siamltex}
\documentclass{report}

%\input{../China2012/mypackage.tex}
%\input{../China2012/mydefinition.tex}

% publish double column
%\documentclass{../elsevier/elsart3}

% manuscript
% Use the option doublespacing or reviewcopy to obtain double line spacing
%\documentclass[doublespacing]{../elsevier/elsart}

%\setlength{\topmargin}{0.2cm}
%\setlength{\headheight}{1.1cm}

% if you use PostScript figures in your article
% use the graphics package for simple commands
% \usepackage{graphics}
% or use the graphicx package for more complicated commands
\usepackage{graphicx}
% or use the epsfig package if you prefer to use the old commands
% \usepackage{epsfig}
\usepackage[titletoc,toc,title]{appendix}
%\addtocontents{toc}{\protect\contentsline{chapter}{Appendix:}{}}
\usepackage[textwidth=20cm,textheight=24cm]{geometry}

% The amssymb package provides various useful mathematical symbols
\usepackage{amssymb}
\usepackage{amsmath}

\newcommand{\ds}{\displaystyle}
\newcommand{\bm}[1]{\mbox{\boldmath $#1$}}
\newcommand{\abs}[1]{\left| #1 \right|}
\newcommand{\fff}{\footnotesize}
\renewcommand{\thesection}{\arabic{section}}
%\renewcommand{\thesubsection}{\arabic{section}.\arabic{subsection}}
\renewcommand{\theequation}{\arabic{equation}}
\newcommand{\ff}[2]{\frac{#1}{#2}}
\newcommand{\pp}[2]{\frac{\partial{#1}}{\partial{#2}}}
\newcommand{\psp}[2]{\frac{\partial^2{#1}}{\partial{#2}^2}}
\newcommand{\ppp}[3]{\frac{\partial^2{#1}}{\partial{#2}\partial{#3}}}
\newcommand{\ptp}[4]{\frac{\partial}{\partial{#1}}{#2}\frac{\partial{#3}}{\partial{#4}}}
\def\A{\mbox{\bf A}}
\def\B{\mbox{\bf B}}
\def\R{\mbox{\bf R}}
\def\l{\mbox{\bf l}}
\def\v{\mbox{\bf v}}
\def\x{\mbox{\bf x}}
\def\U{\mbox{\bf U}}


\title{GTS User Manual For Finite Element Solver}
\author{gts team\\
        Princeton Plasma Physics Lab\\
        P.O. Box 451, Princeton, NJ 08543\\
%        {\tt jchen@pppl.gov}
       }
\date{\today}


\begin{document}
\maketitle
\tableofcontents
%\addcontentsline{toc}{section}{Unnumbered Section}

%N
\section{Numerical integration}
\begin{equation}
\begin{array}{l}
\bigl \langle
P
\bigr \rangle
=
\bigl \langle
\begin{array}{lllllllllllllllllll}
1 
&\xi   &\eta 
&\xi^2 &\xi\eta     &\eta^2 
&\xi^3 &\xi^2\eta   &\xi\eta^2               &\eta^3 
&\xi^4 &\xi^3\eta   &\xi^2\eta^2             &\xi\eta^3              &\eta^4 
&\xi^5-5\xi^3\eta^2 &\xi^2\eta^3-\xi^3\eta^2 &\eta^5-5\xi^2\eta^3
\end{array}
\bigr \rangle
\end{array}
\end{equation}

%N N_xi N_eta N_xixi N_xieta N_etaeta
\begin{equation}
\begin{array}{l}
\bigl \langle
P\\
P_\xi\\
P_\eta\\
P_{\xi\xi}\\
P_{\xi\eta}\\
P_{\eta\eta}\\
\bigr \rangle
= \\
%\bigl \langle
\begin{array}{lllllllllllllllllll}
1 
&\xi   &\eta 
&\xi^2 &\xi\eta     &\eta^2 
&\xi^3 &\xi^2\eta   &\xi\eta^2               &\eta^3 
&\xi^4 &\xi^3\eta   &\xi^2\eta^2             &\xi\eta^3              &\eta^4 
&\xi^5-5\xi^3\eta^2 &\xi^2\eta^3-\xi^3\eta^2 &\eta^5-5\xi^2\eta^3
\\
0 
&1      &0
&2\xi   &\eta         &0
&3\xi^2 &2\xi\eta     &\eta^2                  &0
&4\xi^3 &3\xi^2\eta   &2\xi\eta^2              &\eta^3         &0
&5\xi^4-15\xi^2\eta^2 &2\xi\eta^3-3\xi^2\eta^2 &-10\xi\eta^3
\\
0 
&0            &1
&0            &\xi                     &2\eta
&0            &\xi^2                   &2\xi\eta               &3\eta^2 
&0            &\xi^3                   &2\xi^2\eta             &3\xi\eta^2  &4\eta^3 
&-10\xi^3\eta &3\xi^2\eta^2-2\xi^3\eta &5\eta^4-15\xi^2\eta^2
\\
0 
&0                   &0
&2                   &0                  &0
&6\xi                &2\eta              &0                  &0
&12\xi^2             &6\xi\eta           &2\eta^2            &0              &0
&20\xi^3-30\xi\eta^2 &2\eta^3-6\xi\eta^2 &-10\eta^3
\\
0 
&0      &0
&0      &1            &0
&0      &2\xi         &2\eta                   &0
&0      &3\xi^2       &4\xi\eta                &3\eta^2         &0
&-30\xi^2\eta   &6\xi\eta^2-6\xi^2\eta &-30\xi\eta^2
\\
0 
&0            &0
&0            &0                       &2
&0            &0                       &2\xi                   &6\eta
&0            &0                       &2\xi^2                 &6\xi\eta    &12\eta^2 
&-10\xi^3     &6\xi^2\eta  -2\xi^3     &20\eta^3-30\xi^2\eta
\\
\end{array}
%\bigr \rangle
\end{array}
\end{equation}

%N N_xi N_eta N_xixi N_xieta N_etaeta at node 1 (1,xi=0,eta=0)
%N N_xi N_eta N_xixi N_xieta N_etaeta at node 2 (0,xi=1,eta=0)
%N N_xi N_eta N_xixi N_xieta N_etaeta at node 3 (0,xi=0,eta=1)
\begin{equation}
\begin{array}{l}
\bigl \langle
P\\
P_\xi\\
P_\eta\\
P_{\xi\xi}\\
P_{\xi\eta}\\
P_{\eta\eta}\\
\bigr \rangle
= \\
%\bigl \langle
\begin{array}{l|l|l|l|l|l|l|l|l|l|l|l|l|l|l|l|l|l|l}
 1    &     &     &     &     &     &     &     &     &     &     &     &     &     &     &     &     & \\
\hline
      &1    &     &     &     &     &     &     &     &     &     &     &     &     &     &     &     & \\
\hline
      &     &1    &     &     &     &     &     &     &     &     &     &     &     &     &     &     & \\
\hline
      &     &     &2    &     &     &     &     &     &     &     &     &     &     &     &     &     & \\
\hline
      &     &     &     &1    &     &     &     &     &     &     &     &     &     &     &     &     & \\
\hline
      &     &     &     &     &2    &     &     &     &     &     &     &     &     &     &     &     & \\
\hline
      &     &     &     &     &     &     &     &     &     &     &     &     &     &     &     &     & \\
\hline
 1    &1    &     &1    &     &     &1    &     &     &     &1    &     &     &     &     &1    &     & \\
\hline
      &1    &     &2    &     &     &3    &     &     &     &4    &     &     &     &     &5    &     & \\
\hline
      &     &1    &     &1    &     &     &1    &     &     &     &1    &     &     &     &     &     & \\
\hline
      &     &     &2    &     &     &6    &     &     &     &12   &     &     &     &     &20   &     & \\
\hline
      &     &     &     &1    &     &     &2    &     &     &     &3    &     &     &     &     &     & \\
\hline
      &     &     &     &     &2    &     &     &2    &     &     &     &2    &     &     &-10  &-2   & \\
\hline
      &     &     &     &     &     &     &     &     &     &     &     &     &     &     &     &     & \\
\hline
 1    &     &1    &     &     &1    &     &     &     &1    &     &     &     &     &1    &     &     &1 \\
\hline
      &1    &     &     &1    &     &     &     &1    &     &     &     &     &1    &     &     &     & \\
\hline
      &     &1    &     &     &2    &     &     &     &3    &     &     &     &     &4    &     &     &5 \\
\hline
      &     &     &2    &     &     &     &2    &     &     &     &     &2    &     &     &     &2    &-10 \\
\hline
      &     &     &     &1    &     &     &     &2    &     &     &     &     &3    &     &     &     & \\
\hline
      &     &     &     &     &2    &     &     &     &6    &     &     &     &     &12   &     &     &20 \\
\hline
\\
1 
&\xi   &\eta 
&\xi^2 &\xi\eta     &\eta^2 
&\xi^3 &\xi^2\eta   &\xi\eta^2               &\eta^3 
&\xi^4 &\xi^3\eta   &\xi^2\eta^2             &\xi\eta^3              &\eta^4 
&\xi^5-5\xi^3\eta^2 &\xi^2\eta^3-\xi^3\eta^2 &\eta^5-5\xi^2\eta^3
\\
0 
&1      &0
&2\xi   &\eta         &0
&3\xi^2 &2\xi\eta     &\eta^2                  &0
&4\xi^3 &3\xi^2\eta   &2\xi\eta^2              &\eta^3         &0
&5\xi^4-15\xi^2\eta^2 &2\xi\eta^3-3\xi^2\eta^2 &-10\xi\eta^3
\\
0 
&0            &1
&0            &\xi                     &2\eta
&0            &\xi^2                   &2\xi\eta               &3\eta^2 
&0            &\xi^3                   &2\xi^2\eta             &3\xi\eta^2  &4\eta^3 
&-10\xi^3\eta &3\xi^2\eta^2-2\xi^3\eta &5\eta^4-15\xi^2\eta^2
\\
0 
&0                   &0
&2                   &0                  &0
&6\xi                &2\eta              &0                  &0
&12\xi^2             &6\xi\eta           &2\eta^2            &0              &0
&20\xi^3-30\xi\eta^2 &2\eta^3-6\xi\eta^2 &-10\eta^3
\\
0 
&0      &0
&0      &1            &0
&0      &2\xi         &2\eta                   &0
&0      &3\xi^2       &4\xi\eta                &3\eta^2         &0
&-30\xi^2\eta   &6\xi\eta^2-6\xi^2\eta &-30\xi\eta^2
\\
0 
&0            &0
&0            &0                       &2
&0            &0                       &2\xi                   &6\eta
&0            &0                       &2\xi^2                 &6\xi\eta    &12\eta^2 
&-10\xi^3     &6\xi^2\eta  -2\xi^3     &20\eta^3-30\xi^2\eta
\\
\end{array}
%\bigr \rangle
\end{array}
\end{equation}




% ................................... %
\section{First Of All: How To Get The Code}
\label{sec:svn}
GTS can be obtained from

git clone https://github.com/PrincetonUniversity/PPPL-GTS

\iffalse
A finite element based poisson solver has been implemented in 2013
on triangular mesh of general tokamak geometry with flux coordinates in GTS.
This solver is fully accessible from input parameters from $C^0$ linear element,
$C^0$ 2nd order element, to $C^1$ $5$th order elements.
It can take Dirichlet boundary condtion, Neuman boundary condition,
or mixed Dirichlet-Neuman boundary conditon.
It solves the generated linear systems using direct solver or preconditioned iterative solver.
Derivative finding operator and zonal flow operator are also implemented.

The new finite element Poisson solver has been rigorously tested against analytical solutions
and benchmarked with the previous solver for nonlinear ITG and TEM simulations.
Improved performance for treating nonlinear physics is obtained,
and now it is routinely used for GTS production runs.
Further optimizations for the solver speedup is on the way.
The finite element solver is a key component
for the ongoing development of electromagnetic simulation capability in GTS.
The same finite element structure/framework will be applied to
solve additional equations needed for finite-beta simulations in GTS.
\fi

\iffalse
GTS is under svn version control at

\begin{verbatim}
      svn+ssh://svnsrv.pppl.gov/svn/gts
\end{verbatim}

\noindent
It has a trunk/GTS version and a branches/GTS-BRANCH-FROM-R77 version.
The latter one brached off from trunk at rel77 on June 30, 2017.
A working copy can be checked out from both PPPL and NERSC. But all need you to do

\begin{verbatim}
      module load subverison 
\end{verbatim}

\noindent
to get access. Here is a list of most often used commands:

\begin{verbatim}
1) svn list -R svn+ssh://svnsrv.pppl.gov/svn/gts
\end{verbatim}

\noindent
allows you to check what are in the svn repo.

\begin{verbatim}
2) svn co svn+ssh://svnsrv.pppl.gov/svn/gts/trunk/GTS
   svn co svn+ssh://svnsrv.pppl.gov/svn/gts/branches/GTS-BRANCH-FROM-R77
\end{verbatim}

\noindent
allows you to check out a working copy from PPPL, and

\begin{verbatim}
   svn co svn+ssh://username@svnsrv.pppl.gov/svn/gts/trunk/GTS
   svn co svn+ssh://username@svnsrv.pppl.gov/svn/gts/branches/GTS-BRANCH-FROM-R77
\end{verbatim}

\noindent
allows you to do so from NERSC.

\begin{verbatim}
4) svn status -u
\end{verbatim}

\noindent
allows you to check your working copy status.

\begin{verbatim}
5) svn log 
\end{verbatim}

\noindent
allows you to check commit history.

\begin{verbatim}
6) svn diff
\end{verbatim}

\noindent
allows you to compare the difference of your working copy against the repo.

\begin{verbatim}
7) svn update
\end{verbatim}

\noindent
allows you to update your working copy from repo's newest changes.

\begin{verbatim}
8) svn ci GTS -m "your message"
\end{verbatim}

\noindent
allows you to commit your changes back into the repo  (ex: from GTS parent directory).

\begin{verbatim}
   Note: Before a commit, you must do "svn update".
         After a commit, you must do "svn update" again
         to avoid any future conflicts.
\end{verbatim}

\begin{verbatim}
9) svn update
   svn add filename/dirname
   svn ci GTS -m "your message"
   svn update
\end{verbatim}

\noindent
allows you to add a new file or directory back into GTS repo.
\fi



% ................................... %
\section{Second Of All: How To Compile The Code}
\label{sec:compile}

\begin{verbatim}
make _ARCH=perlmutter
make _ARCH=stellar
make _ARCH=spark
\end{verbatim}


% ................................... %
\section{Third Of All: How To Run The Code}
\label{sec:run}
GTS runs on both EDISON and CORI HASWELL \& KNL at NERSC.
The first input parameter is

\begin{verbatim}
      order 
\end{verbatim}

\noindent
given in the input file with default name

\begin{verbatim}
      INPUT.d
\end{verbatim}

\noindent
and sets the order of linear solver

\begin{verbatim}
      solverobject%order=order
      order = 0 original new_gk_solver solver
      order = 1 new solver with linear finite element
      order = 2 new solver with 2nd order finite element
\end{verbatim}

\noindent
where the $solverobject$
represents globally declared $solver0$, $solver1$, $solver\_dphi\_dt$, or $masssolver$ in 

\begin{verbatim}
      main.F90
\end{verbatim}

\noindent
When

\begin{verbatim}
      order >= 1
\end{verbatim}

\noindent
the $solverobject$ is created through

\begin{verbatim}
      new_fem_solver
\end{verbatim}

\noindent
Otherwise, 

\begin{verbatim}
      new_gk_solver
\end{verbatim}

\noindent
is to be called as the default solver.  
$order$ is also assigned to 

\begin{verbatim}
      solverobject%fem%oder=order
\end{verbatim}

\noindent
to define the order of corresponding finite elements.
For $masssolver$, which finds the partial derivative and Laplacian
in $(a, \theta)$ flux coordinates, we always have

\begin{verbatim}
      masssolver%oder=0
      masssolver%fem%oder=1
\end{verbatim}

\noindent
Imposing boundary condition for $masssolver$ is not necessary.

Each $solverobject$ has its own set of mesh and elements 
so that different set or order of elements can be built from solver to solver 
to give us great flexibility.

In $new\_fem\_solver$, 
a triangular mesh is generated through utility $triangle$,
then finite elements are created on this mesh,
the solver matrix and multiply matrix are created accordingly,
and finally the solve matrix is modified by the type of matrix and boundary conditions.
The code returns the successfully created $solverobject$ 
with its content stored in

\begin{verbatim}
      solverobject%ksp
\end{verbatim}

\noindent
if there is no error.  Not untill the linear system is actually to be solved,
will the right hand side vector $b$ be modified according to section \ref{sec:bc} in subroutine

\begin{verbatim}
      petsc_solve
\end{verbatim}

\noindent
depending on the type of boundary condition.  From

\begin{verbatim}
      INPUT.d
\end{verbatim}

\noindent
the boundaries are given at

\begin{verbatim}
      a0 : a_min is the flux surface closest to magnetic axis 
      a1 : a_max, the outmost flux surface 
\end{verbatim}

\noindent
and can be prescribed with the following $4$ types of boundary conditons

\begin{verbatim}
      bc_type = 0 neumann boundary conditions at both a=a0 and a=a1
      bc_type = 1 dirichlet boundary condition at a=a0,
                  neumann boundary conditions at a=a1
      bc_type = 2 dirichlet boundary conditions at both a=a0 and a=a1
      bc_type = 3 neumann boundary condition at a=a0,
                  dirichlet boundary conditions at a=a1
\end{verbatim}

\noindent
being passed as an argument to new\_fem\_solver.  
In the case of $solver\_dphi\_dt$,
only homogeneous Dirichlet boundary conditions will be imposed on boundaries
at $a=a0$ and $a=a1$.

The last important parameter is 

\begin{verbatim}
      mat_type
\end{verbatim}

\noindent
and is passed to new\_fem\_solver
in the same way as that bc\_type is argumented.

\begin{verbatim}
      mat_type == 2: alpha=0
      mat_type != 2: alpha!=0
\end{verbatim}

\noindent
where $\alpha$(alpha) is the coefficients in equations eq (\ref{eq:poisson2}) and eq (\ref{eq:poisson8})
%mstep: maximum timestep the code will run

 

% ................................... %
\section{GTS Parameters}
\label{sec:para}
\begin{center}
\begin{tabular}{ |l|l| } 
 \hline
  tstep=0.020, &   tstep=0.010, \\ 
  micell=60, &   micell=60, \\ 
  mpsi=160, &   mpsi=160, \\ 
  ipsi=85, &   ipsi=85, \\ 
  mthetamax=1, &   mthetamax=1, \\ 
  mzetamax=32, &   mzetamax=64, \\ 
               &                \\
on edison: npartdom=24, or multiple of 24 &   on cori: npartdom=32, or multiple of 32 \\ 
  ipsi=mpsi/2  & ipsi=mpsi/2\\
& \\
 \hline
\end{tabular}
\end{center}
\\

The number of mpi tasks is determined by the product of mzetamax and npartdom, for example

\begin{verbatim}
      mzetamax=32,
      npartdom=24,
      ntasks=32*24=768
\end{verbatim}

 

% ................................... %
\section{GTS Mesh Grids and Safety Factor}
\label{sec:q}
GTS mesh setup has been described in reference \cite{WW1}.
Here we only record the part important to understand finite element mesh.
The GTS mesh is constructed using modified safety factor $\bar{q}(a)$,
which is used everywhere in the code.
In the magnetic flux coordinates $(a, \theta, \zeta)$, 
the field lines are straight and
the radial coordinate labels magnetic surface
$$
a = \sqrt{\ff{\psi}{\psi_e}}
$$
where $\psi$ and $\psi_e$ are the toroidal flux and its value on the plasma boundary, 
repectively. Therefore,

$$
a0 \leq a \leq a1,\;\;\; 0<a0<a1 \leq 1
$$
$a0$ and $a1$ are input parameters explained in section (\ref{sec:run}).

\iffalse
Two sets of mesh are constructed in this coordinate systems.
One set is 2D mesh with a uniform grid
in both radial and poloidal directions to represent the axisymmetric MHD equilibrium.
This mesh has a larger grid size
since the equilibrium scale length is much larger than that of the turbulence.
Then multidimensional spline interpolation is needed to derive the equilibrium quatities
at any spatial location.

The other set of mesh is called field-line-following mesh.
In this mesh, the symmetriy is broken and 
non-symmetric flux coordinate has to be used.
If we rescale radial coordinate $a$ by defining $\rho$ as
$$
\ff{d\rho}{da} = \sqrt{\ff{T_c}{T_i(a)}}
$$
where $T_c$ is the temperature at a reference point radial location,
then evenly spaced radial grid can be applied to the new radial coordiante $\rho$.
This allows the grid size in real space to be correlated with the local gyro-radius
$$
\Delta a \sim \sqrt{\ff{T_i(a)}{T_c}}
$$
\fi

In the poloidal direction, the total number of radial grids,
each one representing one flux surface, is given by the global parameter

\begin{verbatim}
      mpsi
\end{verbatim}

\noindent
in

\begin{verbatim}
      INPUT.d
\end{verbatim}

\noindent
and the $\theta$ grid with grid size $\Delta \theta(a)$ is computed and saved in array

\begin{verbatim}
      mtheta(0:mpsi)
\end{verbatim}

\noindent
which is uniform on each flux surface, while varying from surface to surface.
It is determined so as to make poloidal arc length
$\Delta l_\theta$ near the midplane correlated with $\rho_i$ in the loop

\begin{verbatim}
      do i=0,mpsi
         tdum= ...
         mtheta(i)= ... tdum ...
      enddo
\end{verbatim}

\noindent
in code

\begin{verbatim}
      setup_v2.F90
\end{verbatim}

\noindent
and the total number of grid on each poloidal plane is evaluated and saved in parameter

\begin{verbatim}
      mgrid=sum(mtheta+1)
\end{verbatim}

\noindent
Note that the mesh points are doubly counted when $\theta=2\pi$ meets $\theta=0$ in $mgrid$
caused by polidal $2\pi$ periodicity.
%A typical mesh in the poloidal plane $\zeta=0$ is set up in figure 1.
By excluding this repeated redundancy, the total number of grids on each poloidal plane is computed by

\begin{verbatim}
      neq=sum(mtheta)
\end{verbatim}

\noindent
where $\left[ 1:mtheta(0) \right]$ corresponds to boundary at $a=a0$, and
      $\left[ neq-mtheta(mpsi-1)+1:neq \right]$ corresponds to boundary at $a=a1$.

Once the 2D mesh is constructed on the $\zeta=0$ plane,
a 3D mesh is constructed by starting from a grid point,
following each (approximate) field line which satisfies
$$
\bar{q}(a)\theta - \zeta = const
$$
$\bar{q}(a)$ is slightly different from the usual safety factor $q(a)$ 
and makes approximate come back to one of the grid points on the $\zeta=0$ plane.

\iffalse
Two methods are used to enforce the toroidal periodicity.
One maps the grid at $\zeta=0$ to the grid at $\zeta=2\pi$ using interpolation;
the other one allows the grid to depart slightly from the magnetic field lines
to match the grid points. 
The first method will result in some spatial damping.
The second method requires a chain rule in calculating the parallel derivatives.

Interpolation and deposition will transfer information between particles and grid points.
They are performed within a "flux-tube type" of cell twisted along the field line 
rather than a toroidally symmetric cell.
\fi

The number of toroidal grids, i.e., the number of poloidal planes,

\begin{verbatim}
      mzetamax
\end{verbatim}

\noindent
is chosen as input parameter to give adequate resolution
for the parallel structure of the modes.

\iffalse
The gyro-kinetic particles are followed in general flux coordinatess 
using guiding center Lagrangian equations. 
The guiding center Lagrangian has the following normalized form:
$$
L(\x, \dot{\x}; t) = (\A + \rho_\parallel \B) \cdot \v - H
$$
$H$ is the guiding center Hamiltonian
$$
H=\rho^2_\parallel B^2/2 + \mu B + \Phi.
$$
$\B = \boldsymbol{\nabla} \times \A$, 
$\rho_\parallel = v_\parallel / B$ is the parallel gyro-radius,
$\mu$ is the magnetic moment,
and $\Phi$ is the electric potential.
The independent variables are 
$\x = \x(r, \theta, \zeta, \rho_\parallel)$.
The Lagrangian equation for the motion is given by
$$
\ff{d}{dt}\left( \pp{}{x_i} L \right) - \pp{}{x_i} L = 0
$$
The obtained equations for $\ff{d\x}{dt}$ are suitable for any generalized flux coordinates.
\fi


% ................................... %
\section{Poisson Equation and new\_gk\_solver}
\label{sec:new_gk_solver}
new\_gk\_solver uses four-point averaging scheme to solve GK poisson eq.

Generally, with distribution function defined as a 7D function 
$f(\x, \v, t)$
of real space $\x$, velocity space $\v$, and time $t$, 
the density can be given as a 3D function of real space $\x$
\begin{equation}
\label{eq:density}
n(\x, t) \equiv \int f(\x, \v, t) d^3v
\end{equation}
Similarly, we can defined  the perturbed ion density and electron density
\begin{equation}
\begin{array}{l}
\delta n_i(\x, t) = \int \delta f_i(\x, \v, t) d^3v
\\
\delta n_e(\x, t) = \int \delta f_e(\x, \v, t) d^3v
\end{array}
\end{equation}
in real space.
The quasi-neutrality becomes 
$$\delta n_i(\x, t) \simeq \delta n_e(\x, t)$$

But in gyro-kinetic studies,
the distribution function is defined at guiding center $\R$, as a 6D function
$f_{gc}(\R, \mu, v_\parallel, t)$
of guiding center position $\R$, 
magnetic moment $\mu$, 
the parallel velocity $v_\parallel$,
and time $t$.
Thus, the transformation
between the particle position $\x$ and the guiding center position $\R$
should be applied.

\noindent
Briefly, we have the guiding center density
$$
n_{gc}(\R, t) \equiv \int f_{gc}(\R, \mu, \v_\parallel, t) \mathcal{J} d\mu d\v_\parallel
$$

\noindent
the potential $\phi$
$$
\bar{\phi}(\R, \mu)=\ff{1}{2\pi} \int \phi(\x) \delta(\x - \R - \bm{\rho}) d\x d\Theta
$$
$$
\tilde{\phi}(\x) = \ff{1}{2\pi} \int \bar{\phi}(\R, \mu) f_{0i}(\R, \mu, v_\parallel)
\delta(\R - \x + \bm{\rho}) d\R d\mu dv_\parallel d\Theta
$$

\noindent
the ion density $\delta n_i$
$$
\delta \bar{n}_i(\x) = \ff{1}{2\pi} \int
\delta f_i(\R, \mu, v_\parallel) \delta(\R - \x + \bm{\rho}) d\R d\mu dv_\parallel d\Theta
$$

\noindent
the electron density $\delta n_e$
$$
\delta n_e(\x) = \ff{1}{2\pi} \int
\delta f_e(\R, \mu, v_\parallel) \delta(\R - \x + \bm{\rho}) d\R d\mu dv_\parallel d\Theta
 = \ff{1}{2\pi} \int
\delta f_e(\x, \mu, v_\parallel) d\mu dv_\parallel
$$

\noindent
where $\bm{\rho}$ is the gyro-radius vector,
$\Theta$ is the gyro-phase,
$f_{0i}$ is assumed to be Maxwellian for ions,
and $\delta f_i$ is the perturbed ion distribution function.
Using 
$$
\delta n_i(\x, t) 
= \delta \bar{n}_i(\x, t) + \delta n_i^{polarization}
= \delta \bar{n}_i(\x, t) - n_0 \ff{q_i}{T_i} (\phi - \tilde{\phi})
$$
and the quasi-neutrality condition, we obtain the most used poission equation
\begin{equation}
\label{eq:1a}
\ff{q_i}{T_i} (\phi - \tilde{\phi})
= \ff{\delta \bar{n}_i(\x, t) }{ n_0 } -
  \ff{\delta n_e(\x, t) }{ n_0 } -
\end{equation}



\iffalse
% ................................... %
\section{Gyro-kinetic transformation and the }

%One of the important elements of the gyro-kinetic formalism is 
The fluctuations, such as  and 
in the two coordinates are connected by the gyro-kinetic transformation

In a gyro-kinetic simulation, the quantities are calculated either in real space
or in Fourier space.

In real space, the transformation is carried out by the four-point averaging scheme.
The exact average is performed on a gyro-plane perpendicular to the magnetic field,
with four points exactly spaced on a gyro-orbit.
Because the grid points on which perturbed potential and perturbed density
are calculated are set up on poloidal planes,
it is much more convenient to perform the gyro-average on ploidal planes 
instead of gyro-planes.
In general geometry, the finite ratio of poloidal to total magnetic field 
$\ff{B_p}{B}$ will separate poloidal plane from gyro-plane. 
The projection of gyro-plane onto poloidal plane becomes an elliptic orbit.

In the direction of 
$(\boldsymbol{\nabla} \zeta \times \B) \times \boldsymbol{\nabla} \zeta$,
the gyro-radius is elongated by a factor of $\ff{1}{\cos\alpha}$
while there is no change in the 
$(\boldsymbol{\nabla} \zeta \times \B)$ direction.
Here
$$
\cos\alpha = \ff{\B \cdot \boldsymbol{\nabla}\zeta }{ B |\nabla\zeta| }
= \ff{\psi^\prime }{ \mathcal{J} B | \nabla \zeta | }
$$
with $\psi^\prime \equiv \ff{d\psi}{dr}$ and 
the Jacobian defined as 
$\mathcal{J} = (\nabla r \cdot \nabla \theta \times \nabla \zeta)^{-1} > 0 $.

The four points used for the gyro-average are chosen to be located on the axes
$(\nabla \zeta \times \B) \times \nabla \zeta$ and
$(\nabla \zeta \times \B)$.
An ion spends approximately the same amount of time on each quarter of the ellipse.
To locate the the positions of these points relative to the guiding center,
the directional derivatives need to be determined first.

The directional derivative is defined as
$$
\ff{du}{dl} \equiv \hat{\l} \cdot \boldsymbol{\nabla} u
$$
for any function $u$. 
$\hat{\l}$ is the unit vector in direction $\l$.

In the 
$\l_1 = (\boldsymbol{\nabla} \zeta \times \B)$
direction
$$
\begin{array}{l}
\ff{dr}{dl_1} = \\
\ff{d\theta}{dl_1} = \\
\end{array}
$$

In the 
$\l_2 = (\boldsymbol{\nabla} \zeta \times \B) \times \boldsymbol{\nabla} \zeta$
direction
$$
\begin{array}{l}
\ff{dr}{dl_2} = \\
\ff{d\theta}{dl_2} = \\
\end{array}
$$

The four points are determined by
$$
r_j = r_0 + \Delta r_j,
\theta_j = \theta_0 + \Delta \theta_j,
\zeta_j = \zeta_0, 
j = 1, 2, 3, 4
$$
with
$$
\begin{array}{l}
\Delta r_{1,2} = \pm \ff{dr}{dl_1} \rho_i,
\Delta \theta_{1,2} = \pm \ff{d\theta}{dl_1} \rho_i,
\\
\Delta r_{3,4} = \pm \ff{dr}{dl_2} \ff{\rho_i}{\cos\alpha},
\Delta \theta_{3,4} = \pm \ff{d\theta}{dl_2} \ff{\rho_i}{\cos\alpha}.
\end{array}
$$

The calculation of the potential 
$\tilde{\phi}(\x)$ in terms of $\phi(\x)$ involves
the double averaging process, shown in Fig.3.
$\tilde{\phi}$ at the center point is the four-point average of
$\bar{\phi}$ on the $X$ points,
and 
$\bar{\phi}$ on the $X$ points are the other four-point average of $\phi$ on the $O$ points.
This has been done in different gyro-radii.

Note that the four-point averaging scheme is accurate for
$k_\perp \rho_i \leq 2$
modes.
To resolve shorter wavelength modes, we may use more points for the averaging process.

% ................................... %
\section{Gyro-kinetic equations}

The gyro-kinetic distribution function is expression as the sum of 
equilibrium distribution $f_0$ and turbulence perturbation $\delta f$
$$
f = f_0 + \delta f.
$$

Since gyro-radius is large for ion, the ion particle position $\x$
is different from its guiding center position $\R$. The equation for its 
distribution function is called gyro-kinetic equation.
And the ion gyro-kinetic equation is solved at guiding center $\R$.

While for electron, its gyro-radius is neglectable compared to ion,
and hence the approximation $\x \simeq \R$ can be made and the equaiton
for its distribution function is called drift kinetic equation.
This equation is solved at guiding center $\R$ as well.

In the electrostatic limit, the ion gyro-kinetic equation for $\delta f_i$ is
\begin{equation}
\label{eq:iongyroeq}
%\ref{eq:iongyroeq}
\ff{d \delta f_i}{dt} =
\end{equation}
$\mu$ and $v_\parallel$ are independent variables.
Once $\delta f_i$ is known, 
$\delta \bar{n}_i(\x)$ can be calculated using eq\ref{eq:nbar}.
Note that ion particle position $\x$ is different from its guiding center $\R$
and therefore, $\delta \bar{n}_i(\x)$ and $\delta n_i(\x)$ are different.

The equation for the equilibrium distribution function $f_0$ is given as follows
\begin{equation}
\label{eq:f0}
\ff{df_0}{dt} =
\end{equation}

There are different ways for electrons. Mainly,
one is the full $\delta f$ method,
the other one is the $\delta h$ method.
In each method, one can solve either the total potential 
or the perturbed turbulence potential.


% ................................... %
\subsection{full $\delta f$ electrons}
The electron dynamics is described by the drift kinetic equation, 
neglecting the finite gyro-radius effect
\begin{equation}
\label{eq:edrifteq_df}
%\ref{eq:edrifteq_df}
\ff{d \delta f_e}{dt} =
\end{equation}
Once $\delta f_e$ is known, 
$\delta n_e(\x) \simeq \delta \bar{n}_e(\x)$ can be calculated similar to eq\ref{eq:nbar}
\fi






\iffalse
When $\delta \bar{n}_i(\x)$ and $\delta n_e(\x)$ are solved, we can form 
the following poission equation 
\begin{equation}
\label{eq:1a}
%\ref{eq:1a}
\ff{q_i}{T_i} (\phi - \tilde{\phi}) = 
 \ff{\delta \bar{n}_i}{n_0} - \ff{\delta n_e}{n_0}
\end{equation}
from $0 \simeq \nabla^2 \phi = q(n_i - n_e)$,
$n_i = \bar{n}_i + n_i^{\mbox{polarization}}$
subject to quasi-neutrality condition.
$\phi$ is the total potential.
\fi


\noindent
Instead of solving total potential, 
one may solve for the turbulence contribution $\delta \phi$ 
and zonal flow contribution $\langle \phi \rangle$ 
given by
$$
\delta \phi = \phi - \langle \phi \rangle
$$
The equation for $\delta \phi$ is derived by
taking flux surface average of the total potential equation (\ref{eq:1a})
and subtracting this average from (\ref{eq:1a})
\begin{equation}
\label{eq:1b}
%\ref{eq:1b}
\ff{q_i}{T_i} (\delta \phi - \delta \tilde{\phi}) = 
 \ff{\delta \bar{n}_i - \langle \delta \bar{n}_i \rangle}{n_0} - 
 \ff{\delta n_e - \langle \delta n_e \rangle}{n_0}
\end{equation}

\iffalse
obviously We need the equation to solve for zonal flow 
$\langle \phi \rangle$:
$$
\ff{d\langle \phi \rangle}{da} = \cdots
$$

% ................................... %
\subsection{$\delta h$ electrons}
\fi


An alternate method to solve for electorns is the so called $\delta h$ method.
\iffalse
The electron guiding center distribution function is decomposed into
$$
f_e = f_{e0} - \ff{e\delta \phi}{T_e} f_{e0} + \delta h_e
$$
Again, $\delta \phi$ is the turbulence potential.
The electron equilibrium dristribution function $f_{e0}$ satisfies
the electron version of eq (\ref{eq:f0})
and can be approximated by a shifted Maxwellian containing a parallel flow,
similiar to that for ions.
The second term for $f_e$ represents the adiabatic electron response to the potential
fluctuation due to the fast electron motion.

The nonadiabatic electron distribution function $\delta h_e$
is determined by
\begin{equation}
\label{eq:edrifteq_dh}
%\ref{eq:edrifteq_dh}
\ff{d \delta h_e}{dt} =
\end{equation}
Once $\delta h_e$ is solved, $\delta n_e$ can be obtained and
potential $\phi$ can solved.
As before, 
\fi
In summary,
one way is to solve the total potential.
\begin{equation}
\label{eq:2a}
%\ref{eq:2a}
\left( 1 + \ff{T_i}{T_e} \right)
\ff{q_i}{T_i} \phi- 
\ff{q_i}{T_i} \tilde{\phi}
= \ff{\delta \bar{n}_i}{n_0} - \ff{\delta n_e^h}{n_0}
  + \ff{q_i}{T_e} \langle \phi \rangle
\end{equation}
and
the other way is to solve equation for the turbulence potential %and zonal flow contribution.
%The poisson equation for turbulence potential is
\begin{equation}
\label{eq:2b}
%\ref{eq:2b}
\left( 1 + \ff{T_i}{T_e} \right)
\ff{q_i }{T_i} \delta \phi - 
\ff{q_i }{T_i} \delta \tilde{\phi}
= 
 \ff{\delta \bar{n}_i - \langle \delta \bar{n}_i \rangle}{n_0} - 
 \ff{\delta n_e^h - \langle \delta n_e^h \rangle}{n_0}
\end{equation}
 

% ................................... %
\section{Poisson Equation and new\_fem\_solver}
\label{sec:new_fem_solver}
new\_fem\_solver uses finite elements to solve partial differential GK equation.

Instead of solving integral equation, 
we can derive and solve the elliptic pde for the total potential $\phi$.
From the total potential eq (\ref{eq:1a}), we obtain
\begin{equation}
\label{eq:poisson1}
-\bm{\nabla}_\perp \cdot \left( \ff{q_i n_0}{B\Omega_i} \bm{\nabla}_\perp \Phi \right) 
= \langle \delta \bar{n}_i \rangle - \delta n_e
\end{equation}
More generally, the following GK poisson equation is implemented in the code
\begin{equation}
\label{eq:poisson2}
\alpha \Phi + \beta \bm{\nabla}_\perp \cdot \left( \sum_s g_s \bm{\nabla}_\perp \Phi \right) = b
\end{equation}
where
$$
\begin{array} {l}
\alpha(a) = \left(\ff{e}{e_0}\right) 
            \left(\ff{m_o}{m_e}\right) 
            \ff{1}{T_e},
\beta(a) = - \ff{1}{n_{eo}(a)},
g_s(a) = \ff{Z_s}{B^2} n_{so} \left(\ff{e_o}{e_s}\right) \left(\ff{m_s}{m_o}\right) 
\\
b = \sum_s \ff{\delta \bar{n}_s}{n_{so}} \left(\ff{n_{so}}{n_{eo}}\right) - \ff{\delta n_e^{na}}{n_{eo}}
    + \left(\ff{e}{e_o}\right) \left(\ff{m_o}{m_e}\right) \ff{\langle\Phi\rangle}{\bar{T}_e},
\langle\Phi\rangle = \ff{\oint \Phi {\cal J} d\varphi d\theta}{\oint {\cal J} d\varphi d\theta}
\\
\delta n_e^{na} \mbox{ is the non-adiabatic electron fluctuation density}
\end{array}
$$
$s$ specifies the ion species. From now on we assume single ion specie $s=1$ for simplicity.
In flux coordinate $(a, \theta, \zeta)$, the second term in eq (\ref{eq:poisson2}) becomes
\begin{equation}
\label{eq:poisson3}
\begin{array} {l}
\bm{\nabla}_\perp \cdot \left( g_s \bm{\nabla}_\perp \Phi \right)  \\
= \bm{\nabla} \cdot \left( g_s \bm{\nabla}_\perp \Phi \right) \\
= 
\ff{1}{\cal J} \pp{}{a} \left[ {\cal J} g_s \bm{\nabla}_\perp \Phi \cdot \bm{\nabla} a \right]
+
\ff{1}{\cal J} \pp{}{\theta} \left[ {\cal J} g_s \bm{\nabla}_\perp \Phi  \cdot \bm{\nabla }\theta \right]
+
\ff{1}{\cal J} \pp{}{\varphi} \left[ {\cal J} g_s \bm{\nabla}_\perp \Phi  \cdot \bm{\nabla} \varphi \right]
\end{array}
\end{equation}
where ${\cal J}$ represents the flux coordinate Jacobian
$\left[ \bm{\nabla} a \times \bm{\nabla} \theta \cdot \bm{\nabla} \varphi \right]^{-1}$.
Please refer the Appendix for a summary of formulae in curved coordinate systems.
The first term in (\ref{eq:poisson3}) can be written as
\begin{equation}
\label{eq:poisson4}
\begin{array} {rl}
\bm{\nabla}_\perp \Phi \cdot \bm{\nabla} a
= &
 \left[
 \bm{\nabla} \Phi - (\hat{b} \cdot \bm{\nabla} \Phi) \hat{b} 
 \right] \cdot \bm{\nabla} a
\\
\simeq &
  \bm{\nabla} \Phi \cdot \bm{\nabla} a
\\
= &
 \left(
 \pp{\Phi}{a} \bm{\nabla} a + \pp{\Phi}{\theta} \bm{\nabla} \theta
 \right) \cdot \bm{\nabla} a
=
 \pp{\Phi}{a} g^{a a} + \pp{\Phi}{\theta} g^{a \theta}
\end{array}
\end{equation}
in an axis-symmetry system.
Similarly, the second term in (\ref{eq:poisson3}) can be written as
\begin{equation}
\label{eq:poisson5}
\begin{array} {rl}
\bm{\nabla}_\perp \Phi \cdot \bm{\nabla} \theta
= &
 \left[
 \bm{\nabla} \Phi - (\hat{b} \cdot \bm{\nabla} \Phi) \hat{b} 
 \right] \cdot \bm{\nabla} \theta
\\
\simeq &
  \bm{\nabla} \Phi \cdot \bm{\nabla} \theta
\\
= &
 \left(
 \pp{\Phi}{a} \bm{\nabla} a + \pp{\Phi}{\theta} \bm{\nabla} \theta
 \right) \cdot \bm{\nabla} \theta
=
 \pp{\Phi}{a} g^{a \theta} + \pp{\Phi}{\theta} g^{\theta \theta}
\end{array}
\end{equation}
and the third term in (\ref{eq:poisson3}) becomes
\begin{equation}
\label{eq:poisson6}
\begin{array}{rl}
\ff{1}{\cal J} \pp{}{\varphi} \left[ {\cal J} g_s \bm{\nabla}_\perp \Phi \cdot \bm{\nabla} \varphi \right]
= & g_s \pp{}{\varphi} \left[ \bm{\nabla}_\perp \Phi \cdot \bm{\nabla} \varphi \right]
=  g_s \pp{}{\varphi} \left[ 
   \bm{\nabla} \Phi - (\hat{b} \cdot \bm{\nabla} \varphi ) \hat{b}
                     \right] \cdot \bm{\nabla} \varphi 
\\
\simeq & g_s \pp{}{\varphi} \left[ 
   \bm{\nabla} \Phi \cdot \bm{\nabla} \varphi 
                     \right]
=  g_s \pp{}{\varphi} \left[ 
   \pp{\Phi}{\varphi} |\bm{\nabla} \varphi |^2
                     \right]
=  g_s |\bm{\nabla} \varphi |^2 \pp{}{\varphi} \pp{\Phi}{\varphi}
\\
= & g_s |\bm{\nabla} \varphi |^2 \pp{}{\varphi} \left[ 
   -\ff{1}{q} \pp{\Phi}{\theta}
                     \right]
=  -\ff{g_s}{q} |\bm{\nabla} \varphi |^2 \pp{}{\varphi} \pp{\Phi}{\theta}
\\
\simeq & \ff{g_s}{q^2} |\bm{\nabla} \varphi |^2 \psp{\Phi}{\theta}
=  \ff{g_s}{q^2} \psp{\Phi}{\theta} g^{\varphi \varphi}
\end{array}
\end{equation}
Subtituting
(\ref{eq:poisson4}) 
(\ref{eq:poisson5}) 
(\ref{eq:poisson6}) 
into
(\ref{eq:poisson2}),
the partial differential equation to be solved becomes
\begin{equation}
\label{eq:poisson7}
\begin{array}{l}
\alpha \Phi + 
\\
\beta \left[
\ff{1}{\cal J} \pp{}{a} {\cal J} g_s \left( 
           \pp{\Phi}{a} g^{a a} + \pp{\Phi}{\theta} g^{a \theta}
                        \right)
+
\ff{1}{\cal J} \pp{}{\theta} {\cal J} g_s \left( 
           \pp{\Phi}{a} g^{a \theta} + \pp{\Phi}{\theta} g^{\theta \theta}
                        \right)
+
           \ff{g_s}{q^2} \psp{\Phi}{\theta} g^{\varphi \varphi}
                    \right] 
\\
= b
\end{array}
\end{equation}
Here $g^{aa}$, $g^{a\theta}$, $g^{\theta \theta}$
represent the coordinate metrics.
This equation is solved in solver1.


\iffalse
% ................................... %
\section{flux-surface averaged GK poisson equation}
\label{sec:poisson2}
\fi




If we decompose $\Phi$ into $2$ parts
$$
\Phi = \langle \Phi \rangle + \delta \Phi
$$
\iffalse
$$
\beta \ff{d}{d\nu} \left[
    \sum_s \langle g_s(a,\theta) g^{aa} \rangle \nu_a^\prime \ff{d\langle \Phi \rangle}{da}
                   \right]
=
- \beta \ff{d}{d\nu} \sum_s 
  \langle g_s(a,\theta) \bm{\nabla}_\perp \delta \Phi \cdot \bm{\nabla}a \rangle \nu_a^\prime
+ \langle b \rangle
$$
which can be convertd into
$$
\ff{d}{d\nu} \left[
    \sum_s \langle g_s(a,\theta) g^{aa} \rangle \nu_a^\prime \ff{d\langle \Phi \rangle}{da}
                   \right]
=
-  \ff{d}{d\nu} \sum_s 
  \langle g_s(a,\theta) \bm{\nabla}_\perp \delta \Phi \cdot \bm{\nabla}a \rangle \nu_a^\prime
+ \ff{\nu_a^\prime}{\beta}\langle b \rangle
$$
With
$$
\bm{\nabla}_\perp \delta \Phi \cdot \bm{\nabla}a
=
\pp{\delta \Phi}{a} g^{aa} + \pp{\delta \Phi}{\theta} g^{a\theta}
$$
we can solve the other form of problem, 
\fi
Similarly to eq (\ref{eq:poisson2}) the equation for $\delta \Phi$ will be
\begin{equation}
\label{eq:poisson8}
\alpha \delta \Phi 
+ \beta 
\bm{\nabla}_\perp \cdot \sum_s g_s \bm{\nabla}_\perp \delta \Phi 
=
b - \beta \bm{\nabla}_\perp \cdot \sum_s g_s \ff{d\langle \Phi \rangle}{da} \bm{\nabla}a
\end{equation}
with an additional term on the rhs.
 This addtional term can be extended as follows
$$
\begin{array}{ll}
- \beta \bm{\nabla}_\perp \cdot \sum_s g_s \ff{d\langle \Phi \rangle}{da} \bm{\nabla}a
& =
- \beta \ff{1}{\cal J} \pp{}{a} \left(
           {\cal J} g_s \ff{d\langle \Phi \rangle}{da} \bm{\nabla}a \cdot \bm{\nabla}a
                                \right)
- \beta \ff{1}{\cal J} \pp{}{\theta} \left(
           {\cal J} g_s \ff{d\langle \Phi \rangle}{da} \bm{\nabla} \theta \cdot \bm{\nabla}a
                                     \right)
\\
& =
- \beta \ff{1}{\cal J} \pp{}{a} \left(
           {\cal J} g_s \ff{d\langle \Phi \rangle}{da} g^{aa}
                                \right)
- \beta \ff{1}{\cal J} \pp{}{\theta} \left(
           {\cal J} g_s \ff{d\langle \Phi \rangle}{da} g^{a\theta}
                                     \right)
\end{array}
$$
Thus eq (\ref{eq:poisson8}) has a similar form as 
eq (\ref{eq:poisson7}) with above additional writing on the right hand side
and is solved in solver0.

$solver\_dphi\_dt$ solves the following eq, which is similar to eq (\ref{eq:poisson2}).
The only differences come from setting $\alpha=0$, the unknows $\pp{\Phi}{t}$ to be solved,
and rhs vector $b$.
\begin{equation}
\label{eq:poisson9}
\beta (a) \bm{\nabla}_\perp \cdot \sum_s g_s \bm{\nabla}_\perp \pp{\Phi}{t} = b,
\;\;\;
b=\sum_i \ff{Z_i}{n_{eo}} \pp{\delta \bar{n}_i}{t} 
  - \ff{1}{n_{eo}} \pp{\delta n_e^h}{t} 
  - \pp{}{t}\left( \ff{e\delta \Phi}{T_e} \right)
\end{equation}
The only boundary condition is Dirichlet type $\pp{\Phi}{t}=0$ at $a=a0$ and $a=a1$.



% ................................... %
\section{Finite Elements Method}
Finite elements are built upon the reference triangle:
$$
(\bm{\xi 1}(0,0), \bm{\xi 2}(1,0), \bm{\xi 3}(0,1))
$$
which is mapped from real triangle:
$$
(\bm{x1}(a1,\theta 1), \bm{x2}(a2,\theta 2), \bm{x3}(a3,\theta 3))
$$
using transformation Jacobian $J$ 
\begin{equation}
J = \left| \begin{array}{ll}
        \pp{a}{\xi}  & \pp{\theta}{\xi} \\
        \pp{a}{\eta} & \pp{\theta}{\eta}
    \end{array} \right|
  = \pp{a}{\xi}  \pp{\theta}{\eta} - \pp{a}{\eta} \pp{\theta}{\xi}
\end{equation}
from real space $\bm{x}(a,\theta)$ to reference space $\bm{\xi}(\xi,\eta)$
$$
\bm{\xi} \Longrightarrow \bm{x} = \bm{x}(\bm{\xi})
$$
The reverse transformation $J^{-1}$ is written as
\begin{equation}
J^{-1} = \left| \begin{array}{ll}
             \pp{\xi}{a} & \pp{\eta}{a} \\
             \pp{\xi}{\theta} & \pp{\eta}{\theta}
         \end{array} \right|
\end{equation}
from reference space $\bm{\xi}(\xi,\eta)$ to real space $\bm{x}(a,\theta)$
$$
\bm{x} \Longrightarrow \bm{\xi} = \bm{\xi}(\bm{x})
$$
Three types of elements are available in the current version of GTS
and are represented by $N_i$ for the $ith$ base function.

Now we introuduce the variational formulation by
multiplying quation (\ref{eq:poisson7}) with base function $N_i (i=1, 2, \cdots)$ on its both side, 
integrating over each triangle, and then assembling them together
\begin{equation}
\label{eq:variation1}
\begin{array}{lll}
\alpha \int\int \Phi N_i d\sigma
& + & \beta \int\int \left[
      \ff{1}{\cal J} \pp{}{a} {\cal J} g_s \left( 
           \pp{\Phi}{a} g^{a a} + \pp{\Phi}{\theta} g^{a \theta}
                                           \right)
                   \right] N_i d\sigma
\\
& + & \beta \int\int \left[
      \ff{1}{\cal J} \pp{}{\theta} {\cal J} g_s \left( 
           \pp{\Phi}{a} g^{a \theta} + \pp{\Phi}{\theta} g^{\theta \theta}
                                                \right)
                   \right] N_i d\sigma
\\
& + & \beta \int\int \left[
      \ff{g_s}{q^2} \psp{\Phi}{\theta} g^{\varphi \varphi}
                   \right] N_i d\sigma
\\
= \int\int b N_i d\sigma, & & i=1, 2, \cdots.
\end{array}
\end{equation}
$d\sigma$ is the 2D differential area in real space.
Using the approximated variable $\Phi$ and rhs vector $b$
$$
\Phi = \sum_j \Phi_j N_j, \;\;\;
b = \sum_j b_j N_j
$$
the first term in the eq (\ref{eq:variation1}) is converted to
\begin{equation}
\label{eq:variation2}
\begin{array}{ll}
\alpha \int\int \Phi N_i d\sigma
& = \alpha \int\int \sum_j \Phi_j N_j N_i {\cal J} da d\theta
\\
& = \alpha \sum_j \left(\int\int N_j N_i {\cal J} dad\theta\right) \Phi_j
\\
& = \alpha \sum_j {\cal J}_j m_{ij} \Phi_j
\end{array}
\end{equation}
Similarly, we have for the right-hand-side $b$ term
\begin{equation}
\label{eq:variation3}
\begin{array}{ll}
\int\int b N_i d\sigma
& = \int\int \sum_j b_j N_j N_i {\cal J} dad\theta
\\
& = \sum_j \left(\int\int N_j N_i {\cal J} da d\theta\right) b_j
\\
& = \sum_j {\cal J}_j m_{ij} b_j
\end{array}
\end{equation}
Where ${\cal J}$ is the flux coordiante Jacobian evaluated at node $j$,
and $m_{ij}$ is the mass matrix relted to $N_i$ and $N_j$.
Taking partial integration, the $2nd$ part on the left of eq (\ref{eq:variation1}) becomes
\begin{equation}
\label{eq:variation4}
\begin{array}{l}
\beta \int\int \ff{1}{\cal J} \pp{}{a} \left[
      {\cal J} g_s \left( 
           \pp{\Phi}{a} g^{a a} + \pp{\Phi}{\theta} g^{a \theta}
                            \right)
                              \right] N_i d\sigma
+\\
\beta \int\int \ff{1}{\cal J} \pp{}{\theta} \left[
      {\cal J} g_s \left( 
           \pp{\Phi}{a} g^{a \theta} + \pp{\Phi}{\theta} g^{\theta \theta}
                                 \right)
                              \right] N_i d\sigma
+\\
\beta \int\int \left[
      \ff{g_s}{q^2} \psp{\Phi}{\theta} g^{\varphi \varphi}
               \right] N_i d\sigma
 \\
=
\beta \int\int \ff{1}{\cal J} \pp{}{a} \left[
      {\cal J} g_s \left( 
           \pp{\Phi}{a} g^{a a} + \pp{\Phi}{\theta} g^{a \theta}
                            \right)
                              \right] N_i  {\cal J} dad\theta
+\\
\;\;\;
\beta \int\int \ff{1}{\cal J} \pp{}{\theta} \left[
      {\cal J} g_s \left( 
           \pp{\Phi}{a} g^{a \theta} + \pp{\Phi}{\theta} g^{\theta \theta}
                                 \right)
                              \right] N_i  {\cal J} dad\theta
+\\
\;\;\;
\beta \int\int \left[
      \ff{g_s}{q^2} \psp{\Phi}{\theta} g^{\varphi \varphi}
               \right] N_i  {\cal J} dad\theta
\\
=
- \beta \int\int 
      \pp{N_i}{a} {\cal J} g_s \left( 
           \pp{\Phi}{a} g^{a a} + \pp{\Phi}{\theta} g^{a \theta}
                               \right) dad\theta
\\
\;\;\;\;
- \beta \int\int 
      \pp{N_i}{\theta} {\cal J} g_s \left( 
           \pp{\Phi}{a} g^{a \theta} + \pp{\Phi}{\theta} g^{\theta \theta}
                                    \right) dad\theta
\\
\;\;\;\;
- \beta \int\int 
      \pp{ }{\theta}\left(
                          {\cal J}\ff{g_s}{q^2} g^{\varphi \varphi} N_i
                    \right)
      \pp{\Phi}{\theta} dad\theta
\\
\simeq
- \beta \int\int \left( {\cal J} g_s g^{a a} \right)
                 \pp{N_i}{a} \pp{\Phi}{a} dad\theta
- \beta \int\int \left( {\cal J} g_s g^{a \theta} \right)
                 \pp{N_i}{a} \pp{\Phi}{\theta} dad\theta
\\
\;\;\;\;
- \beta \int\int \left( {\cal J} g_s g^{a \theta} \right)
                 \pp{N_i}{\theta} \pp{\Phi}{a} dad\theta
- \beta \int\int \left( {\cal J} g_s g^{\theta \theta} \right)
                 \pp{N_i}{\theta} \pp{\Phi}{\theta} dad\theta
\\
\;\;\;\;
- \beta \int\int 
      \left( {\cal J}\ff{g_s}{q^2} g^{\varphi \varphi} \right)
      \pp{N_i}{\theta} \pp{\Phi}{\theta} dad\theta
\\
=
- \beta \sum_j \Phi_j \int\int \left( {\cal J} g_s g^{a a} \right)
                 \pp{N_i}{a} \pp{N_j}{a} dad\theta
- \beta \sum_j \Phi_j \int\int \left( {\cal J} g_s g^{a \theta} \right)
                 \pp{N_i}{a} \pp{N_j}{\theta} dad\theta
\\
\;\;\;\;
- \beta \sum_j \Phi_j \int\int \left( {\cal J} g_s g^{a \theta} \right)
                 \pp{N_i}{\theta} \pp{N_j}{a} dad\theta
- \beta \sum_j \Phi_j \int\int \left( {\cal J} g_s g^{\theta \theta} \right)
                 \pp{N_i}{\theta} \pp{N_j}{\theta} dad\theta
\\
\;\;\;\;
- \beta \sum_j \Phi_j \int\int 
      \left( {\cal J}\ff{g_s}{q^2} g^{\varphi \varphi} \right)
      \pp{N_i}{\theta} \pp{N_j}{\theta} dad\theta
\\
=
- \beta \sum_j \left( {\cal J} g_s g^{a a} \right)_{ij} srr_{ij} \Phi_j 
- \beta \sum_j \left( {\cal J} g_s g^{a \theta} \right)_{ij} srz_{ij} \Phi_j
\\
\;\;\;\;
- \beta \sum_j \left( {\cal J} g_s g^{a \theta} \right)_{ij} szr_{ij} \Phi_j
- \beta \sum_j \left( {\cal J} g_s g^{\theta \theta} \right)_{ij} szz_{ij} \Phi_j
\\
\;\;\;\;
- \beta \sum_j \left( {\cal J}\ff{g_s}{q^2} g^{\varphi \varphi} \right)_{ij} szz_{ij} \Phi_j
\end{array}
\end{equation}
Here $srr_{ij}$, 
     $srz_{ij}$, 
     $szr_{ij}$, and 
     $szz_{ij}$ 
are the stiffness matrices corresponding to operators
$\pp{N_i}{a}$, $\pp{N_i}{\theta}$, and $\pp{N_j}{a}$, $\pp{N_j}{\theta}$.
Summing together, we obtain the following linear systems
\begin{equation}
\begin{array}{ll}
\sum_j & \{ \alpha {\cal J}_j m_{ij} \\
       & - \beta \left( {\cal J} g_s g^{a a} \right)_{ij} srr_{ij} 
         -2\beta \left( {\cal J} g_s g^{a \theta} \right)_{ij} srz_{ij}
         - \beta \left( {\cal J} g_s g^{\theta \theta} \right)_{ij} szz_{ij}
\\
       & - \beta \left( {\cal J}\ff{g_s}{q^2} g^{\varphi \varphi} \right)_{ij} szz_{ij}
       \} \Phi_j 
\\
= & \sum_j {\cal J}_j m_{ij} b_j, \;\;\; i= 1, 2, \cdots
\end{array}
\end{equation}
which is implemented in $solver1$.

A similar form for equation (\ref{eq:poisson8}) in section \ref{sec:new_fem_solver} 
can be given easily.
First the additional term on the rhs is
$$
\begin{array}{l}
%\int\int - \beta \bm{\nabla}_\perp \cdot \sum_s g_s \ff{d\langle \Phi \rangle}{da} \bm{\nabla}a N_i d\sigma =
%\\
  - \beta \int\int \ff{1}{\cal J} \pp{}{a} \left(
           {\cal J} g_s \ff{d\langle \Phi \rangle}{da} g^{aa}
                                           \right) N_i d\sigma
  - \beta \int\int \ff{1}{\cal J} \pp{}{\theta} \left(
           {\cal J} g_s \ff{d\langle \Phi \rangle}{da} g^{a\theta}
                                                \right) N_i d\sigma
\\
= - \beta \int\int \ff{1}{\cal J} \pp{}{a} \left(
           {\cal J} g_s g^{aa} \ff{d\langle \Phi \rangle}{da}
                                           \right) N_i {\cal J} dad\theta
  - \beta \int\int \ff{1}{\cal J} \pp{}{\theta} \left(
           {\cal J} g_s g^{a\theta} \ff{d\langle \Phi \rangle}{da}
                                                \right) N_i {\cal J} dad\theta
\\
\simeq - \beta \sum_j \int\int \left( {\cal J} g_s g^{aa} \ff{d\langle \Phi \rangle}{da} \right)_j
                        \pp{N_j}{a} N_i dad\theta
       - \beta \sum_j \int\int \left( {\cal J} g_s g^{a\theta} \ff{d\langle \Phi \rangle}{da} \right)_j
                       \pp{N_j}{\theta} N_i dad\theta
\\
\\
\left[ = - \beta \sum_j sr_{ij} \left( {\cal J} g_s g^{aa} \ff{d\langle \Phi \rangle}{da} \right)_j
  - \beta  \sum_j sz_{ij} \left( {\cal J} g_s g^{a\theta} \ff{d\langle \Phi \rangle}{da} \right)_j \right]
\\
 (=R_{mat} \left( -\beta g_s g^{aa} \ff{d\langle \Phi \rangle}{da} \right) + 
   Z_{mat} \left( -\beta g_s g^{a\theta} \ff{d\langle \Phi \rangle}{da} \right)
\\
\mbox{in } (a, \theta) \mbox{ coordinate, } {\cal J} \mbox{is included in } R_{mat} , Z_{mat}
\\
\\
=\sum_j \left( -\beta g_s g^{aa} \ff{d\langle \Phi \rangle}{da} \right)_j
         \int\int \ff{1}{J} \sum_k E_{jk} \theta_k N_i {\cal J} J d\xi d\eta 
+
 \sum_j \left( -\beta g_s g^{a\theta} \ff{d\langle \Phi \rangle}{da} \right)_j
         \int\int \left(-\ff{1}{J}\right) \sum_k E_{jk} a_k N_i {\cal J} J d\xi d\eta 
\\
 (=R_{mat} \left( -\beta g_s g^{aa} \ff{d\langle \Phi \rangle}{da} \right) + 
   Z_{mat} \left( -\beta g_s g^{a\theta} \ff{d\langle \Phi \rangle}{da} \right)
\\
=\sum_j \left( -\beta g_s g^{aa} \ff{d\langle \Phi \rangle}{da} \right)_j
         \sum_k \theta_k \int\int E_{jk} N_i {\cal J} d\xi d\eta 
+
 \sum_j \left( -\beta g_s g^{a\theta} \ff{d\langle \Phi \rangle}{da} \right)_j
         \sum_k a_k \int\int E_{jk} N_i {\cal J} d\xi d\eta 
\\
=\sum_j \left[ \left( -\beta g_s g^{aa} \ff{d\langle \Phi \rangle}{da} \right)_j
               +
               \left( -\beta g_s g^{a\theta} \ff{d\langle \Phi \rangle}{da} \right)_j
        \right]
         \sum_k \left( \theta_k - a_k\right) P_{ijk}(i,j,k)
\end{array}
$$
where $sr_{ij}$ and $sz_{ij}$ are the partial stiffness matrices corresponding to
operators $\pp{N_j}{a} N_i$ and $\pp{N_j}{\theta} N_i$, repspectively.
The code can be found in subroutine $add\_hh2bb$.
Thus we have the second type of linear systems to be solved
\begin{equation}
\begin{array}{ll}
\sum_j & \{ \alpha {\cal J}_j m_{ij} \\
       & - \beta \left( {\cal J} g_s g^{a a} \right)_{ij} srr_{ij} 
         -2\beta \left( {\cal J} g_s g^{a \theta} \right)_{ij} srz_{ij}
         - \beta \left( {\cal J} g_s g^{\theta \theta} \right)_{ij} szz_{ij}
\\
       & - \beta \left( {\cal J}\ff{g_s}{q^2} g^{\varphi \varphi} \right)_{ij} szz_{ij}
       \} \Phi_j 
\\
\\
\left[ =\right. & \left. \sum_j \{ m_{ij} {\cal J}_j b_j
              - \beta sr_{ij} \left( {\cal J} g_s g^{aa} \ff{d\langle \Phi \rangle}{da} \right)_j
              - \beta sz_{ij} \left( {\cal J} g_s g^{a\theta} \ff{d\langle \Phi \rangle}{da} \right)_j
           \} \right]
\\
\\
= & \sum_j \{ m_{ij} {\cal J}_j b_j
 - \left[ \left( -\beta g_s g^{aa} \ff{d\langle \Phi \rangle}{da} \right)_j
               +
               \left( -\beta g_s g^{a\theta} \ff{d\langle \Phi \rangle}{da} \right)_j
        \right]
         \sum_k \left( \theta_k - a_k\right) P_{ijk}(i,j,k) \}
%\\
\end{array}
\end{equation}
which is implemented in $solver0$.
This is the first option to evaluate this second term
each time the solver is called. 

The second option, which is currently implemented in add\_hh2bb,
follows from finding derivation directly from subroutin find\_dxdadxdt
in section (\ref{sec:dadt}). So the above equation becomes
\begin{equation}
\begin{array}{ll}
\sum_j & \{ \alpha {\cal J}_j m_{ij} \\
       & - \beta \left( {\cal J} g_s g^{a a} \right)_{ij} srr_{ij} 
         -2\beta \left( {\cal J} g_s g^{a \theta} \right)_{ij} srz_{ij}
         - \beta \left( {\cal J} g_s g^{\theta \theta} \right)_{ij} szz_{ij}
\\
       & - \beta \left( {\cal J}\ff{g_s}{q^2} g^{\varphi \varphi} \right)_{ij} szz_{ij}
       \} \Phi_j 
\\
\\
 = & \sum_j m_{ij} {\cal J}_j 
       \left( 
        b_j 
- \beta \ff{1}{\cal J} \pp{}{a} \left(
           {\cal J} g_s \ff{d\langle \Phi \rangle}{da} g^{aa}
                                \right)_j
- \beta \ff{1}{\cal J} \pp{}{\theta} \left(
           {\cal J} g_s \ff{d\langle \Phi \rangle}{da} g^{a\theta}
                                     \right)_j
       \right) 
\end{array}
\end{equation}



% ................................... %
\subsection{$C^0$ linear elements}
\label{sec:e1}
This set of simplest elements is built upon the three nodes of reference triangle and 
has continuity across the element edges:
\begin{equation}
\begin{array}{l}
N_1 = 1-\xi-\eta\\
N_2=\xi\\
N_3=\eta\\
\end{array}
\end{equation}
variable $u(\bm{x})$ to be found in real space can be represent by
\begin{equation}
u(\bm{x})=\sum_j N_j u_j
\end{equation}
where $u_j$ is the variable $u(\bm{x})$ evaluated at triangle node $j$.
With iso-parameter mapping from real triangle to reference triangle,
the coordinates can be written as
\begin{equation}
\begin{array}{l}
a=\sum_j N_j a_j\\
\theta=\sum_j N_j \theta_j\\
\end{array}
\end{equation}
Here $j$ refers to the $j$th node of each triangle. 
Obvisously, the first order derivative of the base functions are found to be
\begin{equation}
\begin{array}{ll}
\pp{N_1}{\xi} = -1   &  \pp{N_1}{\eta} = -1\\
\pp{N_2}{\xi} = 1    &  \pp{N_2}{\eta} = 0\\
\pp{N_3}{\xi} = 0    &  \pp{N_3}{\eta} = 1
\end{array}
\end{equation}
So the Jacobian can be easily computed as 
$$
\begin{array}{ll}
J & = \sum_i a_i \pp{N_i}{\xi}  \sum_j \theta_j \pp{N_j}{\eta} - 
      \sum_i a_i \pp{N_i}{\eta} \sum_j \theta_j \pp{N_j}{\xi}
\\
  & = \sum_i\sum_j a_i  \theta_j \left(\pp{N_i}{\xi}  \pp{N_j}{\eta} - \pp{N_i}{\eta} \pp{N_j}{\xi}\right)
\\
  & = \sum_i\sum_j a_i  \theta_j E_{ij}
\end{array}
$$
The inverse Jacobian $J^{-1}$ is computed below.
We start from the following fact
$$
J \cdot J^{-1} = \left[
                       \begin{array}{ll}
                          \pp{a}{\xi}  & \pp{\theta}{\xi} \\
                          \pp{a}{\eta} & \pp{\theta}{\eta}
                       \end{array}
                 \right]
\cdot
                 \left[
                       \begin{array}{ll}
                          \pp{\xi}{a} & \pp{\eta}{a} \\
                          \pp{\xi}{\theta} & \pp{\eta}{\theta}
                       \end{array}
                 \right]
=
                 \left[
                       \begin{array}{cc}
                          1 & 0 \\
                          0 & 1
                       \end{array}
                 \right]
= I
$$
which gives the following two linear sysems. The first one is
$$
J \left[ \begin{array}{l}
          \pp{\xi}{x} \\
          \pp{\xi}{y} 
         \end{array} \right]
=
  \left[ \begin{array}{l}
          1  \\
          0
  \end{array} \right]
$$
and the second one is
$$
J \left[ \begin{array}{l}
          \pp{\eta}{x} \\
          \pp{\eta}{y}
  \end{array} \right]
=
  \left[ \begin{array}{l}
          0 \\
          1
  \end{array} \right]
$$
Their solution gives us the four entries in the inverse Jacobian $J$
$$
\begin{array}{l}
\pp{\xi}{a} = \ff{1}{J} \left[\begin{array}{ll}
                              1 & \pp{\theta}{\xi} \\
                              0 & \pp{\theta}{\eta} 
                              \end{array}\right]
            = \ff{1}{J} \pp{\theta}{\eta} ;
\\
\pp{\xi}{\theta} = \ff{1}{J} \left[\begin{array}{ll}
                                   \pp{a}{\xi} & 1  \\
                                   \pp{a}{\eta}& 0   
                                   \end{array}\right]
            =-\ff{1}{J} \pp{a}{\eta} ;
\\

\pp{\eta}{a} = \ff{1}{J} \left[\begin{array}{ll}
                               0 & \pp{\theta}{\xi} \\
                               1 & \pp{\theta}{\eta} 
                               \end{array}\right]
             =-\ff{1}{J} \pp{\theta}{\xi} ;
\\
\pp{\eta}{\theta} = \ff{1}{J} \left[\begin{array}{ll}
                                    \pp{a}{\xi} & 0  \\
                                    \pp{a}{\eta}& 1   
                                    \end{array}\right]
             = \ff{1}{J} \pp{a}{\xi} ;
\\
\end{array}
$$
Thus the differentition of base functions in real space $\bm{x}(a,\theta)$ 
can be obtained in the following way.
$$
\begin{array}{l}
\pp{N_i}{a} = \pp{N_i}{\xi} \pp{\xi}{a} + \pp{N_i}{\eta} \pp{\eta}{a}
            = \pp{N_i}{\xi} \ff{1}{J} \pp{\theta}{\eta} + \pp{N_i}{\eta} (-\ff{1}{J}) \pp{\theta}{\xi}
            = \ff{1}{J} (\pp{N_i}{\xi} \pp{\theta}{\eta} - \pp{N_i}{\eta} \pp{\theta}{\xi})
            = \ff{1}{J} E_{i,j} \theta_j
\\
\pp{N_i}{\theta} = \pp{N_i}{\xi} \pp{\xi}{\theta} + \pp{N_i}{\eta} \pp{\eta}{\theta}
                 = \pp{N_i}{\xi} (-\ff{1}{J}) \pp{a}{\eta} + \pp{N_i}{\eta} \ff{1}{J} \pp{a}{\xi}
                 = - \ff{1}{J} (\pp{N_i}{\xi} \pp{a}{\eta} - \pp{N_i}{\eta} \pp{a}{\xi})
                 = - \ff{1}{J} E_{i,j} a_j
\end{array}
$$
Introducing $nw$ GLL points and weights $w(nw)$ on these points,
all the elemental matrices on each triangle can be written using numerical integration:
$$
m_{ij} = \int\int N_i N_j dad\theta
       = \int\int N_i N_j J d\xi d\eta
       = \sum_{n=1}^{n=nw} N_i(n) N_j(n) J(n) w(n)
$$

$$
srr_{ij} = \int\int \pp{N_i}{a}\pp{N_j}{a} dad\theta  
         = \int\int \pp{N_i}{a}\pp{N_j}{a} J d\xi d\eta
         = \sum_{n=1}^{n=nw} \pp{N_i}{a}(n) \pp{N_j}{a}(n) J(n) w(n)
$$

$$
srz_{ij} = \int\int \pp{N_i}{a}\pp{N_j}{\theta} dad\theta  
         = \int\int \pp{N_i}{a}\pp{N_j}{\theta} J d\xi d\eta
         = \sum_{n=1}^{n=nw} \pp{N_i}{a}(n) \pp{N_j}{\theta}(n) J(n) w(n)
$$

$$
szz_{ij} = \int\int \pp{N_i}{\theta}\pp{N_j}{\theta} dad\theta  
         = \int\int \pp{N_i}{\theta}\pp{N_j}{\theta} J d\xi d\eta
         = \sum_{n=1}^{n=nw} \pp{N_i}{\theta}(n) \pp{N_j}{\theta}(n) J(n) w(n)
$$

$$
sr_{ij} = \int\int N_i \pp{N_j}{a} dad\theta
        = \int\int N_i \pp{N_j}{a} J d\xi d\eta
        = \sum_{n=1}^{n=nw} N_i(n) \pp{N_j}{a}(n) J(n) w(n)
$$

$$
sz_{ij} = \int\int N_i \pp{N_j}{\theta} dad\theta
        = \int\int N_i \pp{N_j}{\theta} J d\xi d\eta
        = \sum_{n=1}^{n=nw} N_i(n) \pp{N_j}{\theta}(n) J(n) w(n)
$$



% ................................... %
\subsection{$C^0$ 2nd order elements}
\label{sec:e2}
This set of elements is built upon the linear elements given in the last section
on $6$ nodes (3 nodes and 3 edge nodes) of each triangle.
All the formation of operators and matrices for the 2nd order elements
are exactly the same as section in \ref{sec:e1}.
So we'll just give the base functions and the corresponding first order
derivatives in the reference space here.
\begin{equation}
\begin{array}{l}
N_1 = - \lambda (1 - 2\lambda) \\
N_2 = 4 \lambda \xi \\
N_3 = - \xi ( 1 - 2\xi) \\
N_4 = 4 \xi \eta \\
N_5 = - \eta (1 - 2\eta) \\
N_6 = 4 \eta \lambda
\end{array}
\end{equation}
where $\lambda=1-\xi-\eta$ is the $N_1$, $\xi$ is the $N_2$,
and $\eta$ is the $N_3$ base functions in section \ref{sec:e1},
respectively.
The derivative corresponding to $\xi$ and $\eta$ are
\begin{equation}
\begin{array}{ll}
\pp{N_1}{\xi} = 1-4\lambda   &  \pp{N_1}{\eta} = 1-4\lambda \\
\pp{N_2}{\xi} = 4(\lambda-\xi)    &  \pp{N_2}{\eta} = -4\xi \\
\pp{N_3}{\xi} = -1+4\xi    &  \pp{N_3}{\eta} = 0 \\
\pp{N_4}{\xi} = 4\eta    &  \pp{N_4}{\eta} = 4\xi \\
\pp{N_5}{\xi} = 0    &  \pp{N_5}{\eta} = -1+4\eta \\
\pp{N_6}{\xi} = -4\eta    &  \pp{N_6}{\eta} = 4(\lambda-\eta) \\
\end{array}
\end{equation}
The iso-parameter mapping is used in tranforming from reference space
to real space this set of elements.



% ................................... %
\subsection{$C^1$ 5th order elements}
\label{sec:e5}
This set of elements is much more complicated than the two sets of elements 
introuduced in the previous 2 subsections. Here we have $18$ base functions
built upon the three nodes of each reference triangle,
and each node is related to $6$ of them involving 
$\phi$, 
$\pp{\phi}{a}$,
$\pp{\phi}{\theta}$,
$\psp{\phi}{a}$,
$\ppp{\phi}{a}{\theta}$,
and $\psp{\phi}{\theta}$.
These $18$ base functions are
\begin{equation}
\begin{array}{l}
N_1 = \lambda^2(10\lambda - 15\lambda^2 + 6\lambda^3 + 30\xi\eta(\xi + \eta)) \\
N_2 = \xi \lambda^2(3 - 2\lambda - 3\xi^2  + 6\xi\eta) \\
N_3 = \eta\lambda^2(3 - 2\lambda - 3\eta^2 + 6\xi\eta) \\
N_4 = \ff{1}{2} \xi^2\lambda^2(1 - \lambda + 2\eta) \\
N_5 = \xi^2\eta\lambda^2 \\
N_6 = \ff{1}{2} \eta^2\lambda^2(1 + 2\xi - \eta) \\
\\
N_7 = \xi^2(10\xi^ - 15\xi^2 + 6\xi^3 + 15\eta^2\lambda) \\
N_8 = \ff{1}{2}\xi^2(-8\xi^ + 14\xi^2 - 6\xi^3 - 15\eta^2\lambda) \\
N_9 = \ff{1}{2}\xi^2\eta(6 - 4\xi - 3\eta - 3\eta^2 + 3\xi\eta) \\
N_{10} = \ff{1}{4}\xi^2(2\xi(1 - \xi)^2 + 5\eta^2\lambda) \\
N_{11} = \ff{1}{2}\xi^2\eta(-2 +2\xi + \eta + \eta^2 - \xi\lambda) \\
N_{12} = \ff{1}{4}\xi^2\eta^2\lambda + \ff{1}{2}\xi^3\eta^2\\
\\
N_{13} = \eta^2(10\eta - 15\eta^2 + 6\eta^3 + 15\xi^2\lambda)\\
N_{14} = \ff{1}{2}\xi\eta^2(6 - 3\xi - 4\eta - 3\xi^3 + 3\xi\eta)\\
N_{15} = \ff{1}{2}\eta^2(- 8\eta + 14\eta^2 - 6\eta^3 - 15\xi^2\lambda)\\
N_{16} = \ff{1}{4}\xi^2\eta^2\lambda + \ff{1}{2}\xi^2\eta^3)\\
N_{17} = \ff{1}{2}\xi\eta^2(-2 + \xi + 2\eta + \xi^2 -\xi\eta)\\
N_{18} = \ff{1}{4}\eta^2(2\eta(1-\eta)^2 + 5\xi^2\lambda)\\
\end{array}
\end{equation}
which preseve first order derivative continuity across the element edges. 
Their derivatives in the reference space
$$
\pp{N_i}{\xi}, \;\;\;
\pp{N_i}{\eta}, \;\;\; i=1, 2, \cdots, 18
$$
are simple algebra and will be omitted here.
Check the code for those who are interested in.
The most important numerical issue here is coordinate mapping.
Instead of using iso-parameter mapping as we did in section~\ref{sec:e1}
and section~\ref{sec:e2},
we use sub-parameter mapping 
given by $C^0$ linear element in the section~\ref{sec:e1}, i.e.,
\begin{equation}
\begin{array}{l}
a=\sum_j \bar{N}_j a_j\\
\theta=\sum_j \bar{N}_j \theta_j\\
\end{array}
\end{equation}
to save cost and reduce complexity.
Here the base function $\bar{N}_i$ represents the set of elements given in section~\ref{sec:e1}
and we renam as $\bar{N}_j$ to differentiate it from the $C^1$ elements $N_i$ given in this subsection.
If the precision needs to be increased, the $\bar{N}_i$ can be replaced 
by the second order elements in secton \ref{sec:e2}.
Thus the transformation Jacobian won't change from the previous subsections, i.e.,
$$
\begin{array}{rl}
J = & \pp{a}{\xi}  \pp{\theta}{\eta} - \pp{a}{\eta} \pp{\theta}{\xi} \\
  = & \sum_i a_i \pp{\bar{N}_i}{\xi}  \sum_j \theta_j \pp{\bar{N}_j}{\eta} - 
      \sum_i a_i \pp{\bar{N}_i}{\eta} \sum_j \theta_j \pp{\bar{N}_j}{\xi} \\
  = & \sum_{i,j} (\pp{\bar{N}_i}{\xi}  \pp{\bar{N}_j}{\eta} - 
                  \pp{\bar{N}_i}{\eta} \pp{\bar{N}_j}{\xi}) a_i \theta_j \\
  = & \sum_{i,j} \bar{E}_{i,j} a_i \theta_j
\end{array}
$$
Note the bar sign over the above  opetator $\bar{E}_ij$.
We need to differentiate the base functions in real space $\bm{x}(a,\theta)$
as we did previously:
$$
\begin{array}{ll}
\pp{N_i}{a} & = \pp{N_i}{\xi} \pp{\xi}{a} + \pp{N_i}{\eta} \pp{\eta}{a}
            = \pp{N_i}{\xi} \ff{1}{J} \pp{y}{\eta} + \pp{N_i}{\eta} (-\ff{1}{J}) \pp{y}{\xi}
            = \ff{1}{J} (\pp{N_i}{\xi} \sum_jy_j\pp{\bar{N}_j}{\eta} - \pp{N_i}{\eta} \sum_jy_j\pp{\bar{N}_j}{\xi})
\\
            & = \ff{1}{J} E_{ij} \theta_j
\\
\pp{N_i}{\theta} & = \pp{N_i}{\xi} \pp{\xi}{\theta} + \pp{N_i}{\eta} \pp{\eta}{\theta}
                 = \pp{N_i}{\xi} (-\ff{1}{J}) \pp{a}{\eta} + \pp{N_i}{\eta} \ff{1}{J} \pp{a}{\xi}
                 = - \ff{1}{J} (\pp{N_i}{\xi} \sum_jx_j\pp{\bar{N}_j}{\eta} - \pp{N_i}{\eta} \sum_jx_j\pp{\bar{N}_j}{\xi})
\\
                 & = - \ff{1}{J} E_{ij} a_j
\end{array}
$$
Apparently, $E_{ij}$ here is a mixed product of derivatives of $N_i$ and $\bar{N}_j$,
thus is different from the $E_{ij}$ in sections \ref{sec:e1} and \ref{sec:e2}.
However, the mass and stiffness matrices are computed in the same way as we did in \ref{sec:e1}
with $i,j=1,\cdots,18$.
We rename them with a bar sign over each matrix just as we did for the $\bar{E}_{ij}$ operator:
$$
\bar{m}_{ij}, \;\;\;
\bar{srr}_{ij}, \;\;\;
\bar{srz}_{ij}, \;\;\;
\bar{szz}_{ij}, \;\;\;
\bar{sr}_{ij}, \;\;\;
\bar{sz}_{ij}
$$
Besides this,
the follwoing transformation matrix is needed to convert the higher order derivative from
reference space to real space before the matrices are assembled. 
This is not needed in $C^0$ elements when only $\phi(\bm{x})$ itself is solved.
If you insist to have one, then the transformation matrix in such case would be 
a matrix with one entry $T=[ 1 ]$.

First of all, a new variable $U(\bm{x})$ in real space is formed as
\begin{equation}
U(\bm{x})=\sum_{j=1}^{18} N_j U_j,
\end{equation}
where
\begin{equation}
\begin{array}{lll}
U(\bm{x})= ( & \underbrace{ \phi(\bm{x}_1),
                            \pp{\phi}{a}(\bm{x}_1),
                            \pp{\phi}{\theta}(\bm{x}_1),
                            \psp{\phi}{a}(\bm{x}_1),
                            \ppp{\phi}{a}{\theta}(\bm{x}_1),
                            \psp{\phi}{\theta}(\bm{x}_1)
                 }_{node 1},
& \\
             & \underbrace{ \phi(\bm{x}_2),
                            \pp{\phi}{a}(\bm{x}_2),
                            \pp{\phi}{\theta}(\bm{x}_2),
                            \psp{\phi}{a}(\bm{x}_2),
                            \ppp{\phi}{a}{\theta}(\bm{x}_2),
                            \psp{\phi}{\theta}(\bm{x}_2)
                 }_{node 2},
& \\
             & \underbrace{ \phi(\bm{x}_3),
                            \pp{\phi}{a}(\bm{x}_3),
                            \pp{\phi}{\theta}(\bm{x}_3),
                            \psp{\phi}{a}(\bm{x}_3),
                            \ppp{\phi}{a}{\theta}(\bm{x}_3),
                            \psp{\phi}{\theta}(\bm{x}_3)
                 }_{node 3}
&)^T
\end{array}
\end{equation}
and
$U_j$ is $U(\x)$ evaluated at $\x = (a,\theta)$ on each node.
Our problem is converted to solve for $U(\x)$. Once $U(\x)$ is solved,
we not only have $\phi(\x)$ but also we have its first order derivatives
and second order derivatives. In the reference space $U({\bm \xi})$
has a similar form
\begin{equation}
\begin{array}{lll}
U(\bm{\xi})=( & \underbrace{ \phi(\bm{\xi}_1),
                             \pp{\phi}{\xi}(\bm{\xi}_1),
                             \pp{\phi}{\eta}(\bm{\xi}_1),
                             \psp{\phi}{\xi}(\bm{\xi}_1),
                             \ppp{\phi}{\xi}{\eta}(\bm{\xi}_1),
                             \psp{\phi}{\eta}(\bm{\xi}_1)
                 }_{node 1},
& \\
            &    \underbrace{ \phi(\bm{\xi}_2),
                              \pp{\phi}{\xi}(\bm{\xi}_2),
                              \pp{\phi}{\eta}(\bm{\xi}_2),
                              \psp{\phi}{\xi}(\bm{\xi}_2),
                              \ppp{\phi}{\xi}{\eta}(\bm{\xi}_2),
                              \psp{\phi}{\eta}(\bm{\xi}_2)
                 }_{node 2},
& \\
            &    \underbrace{ \phi(\bm{\xi_3}),
                              \pp{\phi}{\xi}(\bm{\xi}_3),
                              \pp{\phi}{\eta}(\bm{\xi}_3),
                              \psp{\phi}{\xi}(\bm{\xi}_3),
                              \ppp{\phi}{\xi}{\eta}(\bm{\xi}_3),
                              \psp{\phi}{\eta}(\bm{\xi})_3
                 }_{node 3},
& )^T
\end{array}
\end{equation}
From $U(\x)$ to $U(\bm{\xi})$, the transformation matrix at each node $i$ is given as:
$$
T_i =
\left|
\begin{array}{llllll}
1 &        &        &              &                             &              \\
  & J_{11} & J_{12} &              &                             &              \\
  & J_{21} & J_{22} &              &                             &              \\
  &        &        & J_{11}^2     & 2J_{11}J_{12}               & J_{12}^2     \\
  &        &        & J_{11}J_{21} & J_{12}J_{21} + J_{11}J_{22} & J_{12}J_{22} \\
  &        &        & J_{21}^2     & 2J_{21}J_{22}               & J_{22}^2     \\
\end{array}
\right|
$$
It is derived in the following steps.
\begin{equation}
\begin{array}{ll}
\pp{u}{\xi} & = \pp{u}{a}\pp{a}{\xi} + \pp{u}{\theta}\pp{\theta}{\xi}
            = \pp{u}{a} J_{11} + \pp{u}{\theta} J_{12}
\\
\pp{u}{\eta} & = \pp{u}{a}\pp{a}{\eta} + \pp{u}{\theta}\pp{\theta}{\eta}
            = \pp{u}{a} J_{21} + \pp{u}{\theta} J_{22}
\\
\psp{u}{\xi} & = \pp{}{\xi}(\pp{u}{a} J_{11} + \pp{u}{\theta} J_{12})
             = \pp{}{\xi}J_{11}\pp{u}{a} + J_{11}\pp{}{\xi}\pp{u}{a} +
               \pp{}{\xi}J_{12}\pp{u}{\theta} + J_{12}\pp{}{\xi}\pp{u}{\theta}
\\
             & = \pp{}{\xi}J_{11}\pp{u}{a} + \pp{}{\xi}J_{12}\pp{u}{\theta}
               + J_{11}(\pp{}{a}\pp{u}{a}\pp{a}{\xi} + \pp{}{\theta}\pp{u}{a}\pp{\theta}{\xi})
               + J_{12}(\pp{}{a}\pp{u}{\theta}\pp{a}{\xi} + \pp{}{\theta}\pp{u}{\theta}\pp{\theta}{\xi})
\\
             & = \pp{}{\xi}J_{11}\pp{u}{a} + \pp{}{\xi}J_{12}\pp{u}{\theta}
               + J_{11}^2\psp{u}{a} + J_{11}J_{12}\ppp{u}{\theta}{a}
               + J_{12}J_{11}\ppp{u}{a}{\theta} + J_{12}^2\psp{u}{\theta}
\\
\ppp{u}{\xi}{\eta} & = \pp{}{\eta}(\pp{u}{a} J_{11} + \pp{u}{\theta} J_{12})
             = \pp{}{\eta}J_{11}\pp{u}{a} + J_{11}\pp{}{\eta}\pp{u}{a} +
               \pp{}{\eta}J_{12}\pp{u}{\theta} + J_{12}\pp{}{\eta}\pp{u}{\theta}
\\
             & = \pp{}{\eta}J_{11}\pp{u}{a} + \pp{}{\eta}J_{12}\pp{u}{\theta}
               + J_{11}(\pp{}{a}\pp{u}{a}\pp{a}{\eta} + \pp{}{\theta}\pp{u}{a}\pp{\theta}{\eta})
               + J_{12}(\pp{}{a}\pp{u}{\theta}\pp{a}{\eta} + \pp{}{\theta}\pp{u}{\theta}\pp{\theta}{\eta})
\\
             & = \pp{}{\eta}J_{11}\pp{u}{a} + \pp{}{\eta}J_{12}\pp{u}{\theta}
               + J_{11}J_{21}\psp{u}{a} + J_{11}J_{22}\ppp{u}{a}{\theta}
               + J_{12}J_{21}\ppp{u}{\theta}{a} + J_{12}J_{22}\psp{u}{\theta}
\\
\ppp{u}{\eta}{\xi} & = \pp{}{\xi}(\pp{u}{a} J_{21} + \pp{u}{\theta} J_{22})
                   = \pp{}{\xi}J_{21}\pp{u}{a} + J_{21}\pp{}{\xi}\pp{u}{a} +
                     \pp{}{\xi}J_{22}\pp{u}{\theta} + J_{22}\pp{}{\xi}\pp{u}{\theta}
\\
                   & = \pp{}{\xi}J_{21}\pp{u}{a} + \pp{}{\xi}J_{22}\pp{u}{\theta} 
                     + J_{21}(\pp{}{a}\pp{u}{a}\pp{a}{\xi}+\pp{}{\theta}\pp{u}{a}\pp{\theta}{\xi})
                     + J_{22}(\pp{}{a}\pp{u}{\theta}\pp{a}{\xi}+\pp{}{\theta}\pp{u}{\theta}\pp{\theta}{\xi})
\\
                   & = \pp{}{\xi}J_{21}\pp{u}{a} + \pp{}{\xi}J_{22}\pp{u}{\theta} 
                     + J_{21}J_{11}\psp{u}{a}+J_{21}J_{12}\ppp{u}{a}{\theta}
                     + J_{22}J_{11}\ppp{u}{\theta}{a}+J_{22}J_{12}\psp{u}{\theta}
\\
\psp{u}{\eta} & = \pp{}{\eta}(\pp{u}{a} J_{21} + \pp{u}{\theta} J_{22})
              = \pp{}{\eta}J_{21}\pp{u}{a} + J_{21} \pp{}{\eta}\pp{u}{a} +
                \pp{}{\eta}J_{22}\pp{u}{\theta} + J_{22} \pp{}{\eta}\pp{u}{\theta}
\\
              & = \pp{}{\eta}J_{21}\pp{u}{a} + \pp{}{\eta}J_{22}\pp{u}{\theta} 
                + J_{21} (\pp{}{a}\pp{u}{a}\pp{a}{\eta}+\pp{}{\theta}\pp{u}{a}\pp{\theta}{\eta})
                + J_{22} (\pp{}{a}\pp{u}{\theta}\pp{a}{\eta}+\pp{}{\theta}\pp{u}{\theta}\pp{\theta}{\eta})
\\
              & = \pp{}{\eta}J_{21}\pp{u}{a} + \pp{}{\eta}J_{22}\pp{u}{\theta} 
                + J_{21}^2 \psp{u}{a} + J_{21}J_{22}\ppp{u}{a}{\theta}
                + J_{22}J_{21} \ppp{u}{\theta}{a} + J_{22}^2\psp{u}{\theta}
\end{array}
\end{equation}
Thus on each triangle
$$
U(\bm{\xi}) = \bm{T} U(\bm{x}),
\; \; \;
\bm{T}= \left[ \begin{array}{lll}
                  T_1 &    &     \\
                      &T_2 &     \\
                      &    & T_3 \\
               \end{array}
        \right]^T 
$$

$$
mij = \bm{T}^T \cdot \bar{m}_{ij} \cdot \bm{T}
$$

$$
srr_{ij} = \bm{T}^T \cdot \bar{srr}_{ij} \cdot \bm{T},
\; \; \;
srz_{ij} = \bm{T}^T \cdot \bar{srz}_{ij} \cdot \bm{T},
\; \; \;
szz_{ij} = \bm{T}^T \cdot \bar{szz}_{ij} \cdot \bm{T}
$$

$$
sr_{ij} = \bm{T}^T \cdot \bar{sr}_{ij} \cdot \bm{T},
\; \; \;
sz_{ij} = \bm{T}^T \cdot \bar{sz}_{ij} \cdot \bm{T}
$$



% ................................... %
\subsection{Find Derivatives and Laplacian}
\label{sec:dadt}
Given a function $f$, we find derivatives in real space 
$\pp{f}{a}$ and $\pp{f}{\theta}$ 
using the fact
$$
\pp{f}{a} = \pp{f}{a}
$$
Multipying it with base function $N_i$ on both side of the above equaiton
and integrating it over each triangle element
$$
\int\int \pp{f}{a} N_i d\sigma = \int\int \pp{f}{a} N_i d\sigma
$$
i.e.,
$$
\int\int \pp{f}{a} N_i {\cal J} dad\theta = \int\int \pp{f}{a} N_i {\cal J}dad\theta,
$$
in flux coordinate, we have  
\begin{equation}
\label{eq:dfda}
\int\int \sum_j \left(\pp{f}{a}\right)_j N_j N_i {\cal J} J d\xi d\eta = 
\int\int \sum_j f_j \pp{N_j}{a} N_i {\cal J} J d\xi d\eta
\end{equation}
where
$$
f=\sum_j f_j N_j,
\;\;\;
\pp{f}{a} = \sum_j \left(\pp{f}{a}\right)_j N_j
$$
The derivative $\pp{f}{a}$ is found by solving the linear system
$$
\sum_j {\cal J}_j m_{ij} \left(\pp{f}{a}\right)_j
= \sum_j {\cal J}_j sr_{ij} f_j
$$
with the use of $m_{ij}$ and $sr_{ij}$ matrices. 
Or in compact form
$$
M \left(\pp{f}{a}\right) = R_{mat} f
$$


The other way to compute the rhs vector of eq (\ref{eq:dfda}) makes use of operator $P_{ijk}$
\begin{equation}
\label{eq:aPijk}
\begin{array}{ll}
\int\int \sum_j f_j \pp{N_j}{a} N_i {\cal J} J d\xi d\eta
& = \sum_j f_j \int\int \ff{1}{J} \sum_k E_{jk} \theta_k N_i {\cal J} dad\theta
\\
& = \sum_j f_j \sum_k \theta_k \int\int \ff{1}{J} E_{jk} N_i {\cal J} J da\xi d\eta \\
& = \sum_j \sum_k f_j \theta_k P_{ijk} {\cal J}
\end{array}
\end{equation}
In the same way, we can find its $\theta$ derivative by solving the linear system
\begin{equation}
\label{eq:dfdt}
\int\int \sum_j \left(\pp{f}{\theta}\right)_j N_j N_i {\cal J} J d\xi d\eta = 
\int\int \sum_j f_j \pp{N_j}{\theta} N_i {\cal J} J d\xi d\eta
\end{equation}
i.e.,
$$
\sum_j {\cal J}_j m_{ij} \left(\pp{f}{\theta}\right)_j
= \sum_j {\cal J}_j sz_{ij} f_j
$$
In compact form
$$
M \left(\pp{f}{a}\right) = Z_{mat} f
$$

or compute the rhs vetor using $P_{ijk}$
\begin{equation}
\label{eq:tPijk}
\begin{array}{ll}
\int\int \sum_j f_j \pp{N_j}{\theta} N_i {\cal J} J d\xi d\eta
& = - \sum_j f_j \int\int \ff{1}{J} \sum_k E_{jk} a_k N_i {\cal J} dad\theta
\\
& = - \sum_j f_j \sum_k a_k \int\int \ff{1}{J} E_{jk} N_i {\cal J} J da\xi \eta \\
& = - \sum_j \sum_k f_j a_k P_{ijk} {\cal J}
\end{array}
\end{equation}
The computing of (\ref{eq:aPijk}) and (\ref{eq:tPijk}) are implemented in subroutine $find\_dxdadxdt$.
If you just need (\ref{eq:aPijk}) call subroutine $find\_dxda$;
If you just need (\ref{eq:tPijk}) call subroutine $find\_dxdt$.
$find\_dxdt$ calculates the derivative along $\theta$ direction.

According to the term in the square braket right after $\beta$ in eq (\ref{eq:poisson7}):
$$
\left[
\ff{1}{\cal J} \pp{}{a} {\cal J} g_s \left( 
           \pp{\Phi}{a} g^{a a} + \pp{\Phi}{\theta} g^{a \theta}
                        \right)
+
\ff{1}{\cal J} \pp{}{\theta} {\cal J} g_s \left( 
           \pp{\Phi}{a} g^{a \theta} + \pp{\Phi}{\theta} g^{\theta \theta}
                        \right)
+
           \ff{g_s}{q^2} \psp{\Phi}{\theta} g^{\varphi \varphi}
                    \right] 
$$
the Lapacian is implemented in subroutine $find\_laplacian$ using $find\_dxdadxdt$.
%But instead of evaluating derivatives twice to get Laplacian in this way, in section \ref{sec:del} 
%we can derive this quantity by simulating the evalution of first derivatives in this section.



% ................................... %
\subsection{New: Finding Derivatives in $(R,Z)$ Coordinates}
\label{sec:dadt_new}
Considering the accuray, the drivatives found in last section \ref{sec:dadt} can be improved in the following ways.
This time the masssolver is construced in $(R,Z)$ coordinate systems and 
$find\_dxdadxdt$, $find\_dxda$, and $find\_dxdt$
are called to calculate the derivative in $(R,Z)$ coordinate. The the corresponding derivatives
in $(a,\theta)$ coordinate systems can be recovered by

\begin{equation}
\label{tranformation:at2rz}
   \left\{
   \begin{array}{l}
     \pp{\Phi}{a} = \pp{\Phi}{R}\pp{R}{a} + \pp{\Phi}{Z}\pp{Z}{a} \\
     \pp{\Phi}{\theta} = \pp{\Phi}{R}\pp{R}{\theta} + \pp{\Phi}{Z}\pp{Z}{\theta} \\
   \end{array}
   \right.
\end{equation}



% ................................... %
%\subsection{Find 2nd Order Derivative}
%\label{sec:del}
%Given a function $f$, we find its Laplacian in real space 
%$\psp{f}{a}$ and $\psp{f}{\theta}$ 
%using the fact
%$$
%\psp{f}{a} = \psp{f}{a}
%$$
%Multipying it with base function $N_i$ on both side of the above equaiton
%and integrating it over each triangle element
%$$
%\int\int \psp{f}{a} N_i d\sigma = \int\int \psp{f}{a} N_i d\sigma
%$$
%i.e.,
%$$
%\int\int \psp{f}{a} N_i {\cal J} dad\theta = \int\int \psp{f}{a} N_i {\cal J}dad\theta,
%$$
%in flux coordinate, we have  
%\begin{equation}
%\label{eq:dfda}
%\int\int \sum_j \left(\psp{f}{a}\right)_j N_j N_i {\cal J} J d\xi d\eta = 
%\int\int \sum_j f_j \psp{N_j}{a} N_i {\cal J} J d\xi d\eta
%\end{equation}
%where
%$$
%f=\sum_j f_j N_j,
%\;\;\;
%\psp{f}{a} = \sum_j \left(\psp{f}{a}\right)_j N_j
%$$
%The derivative $\psp{f}{a}$ is found by solving the linear system
%$$
%\sum_j {\cal J}_j m_{ij} \left(\psp{f}{a}\right)_j
%= \sum_j {\cal J}_j sr_{ij} f_j
%$$
%with the use of $m_{ij}$ and $sr_{ij}$ matrices.
%Or in compact form
%$$
%M \left(\pp{f}{a}\right) = R_{mat} f
%$$



% ................................... %
\subsection{New: Finite Element In $(R,Z)$ Coordinates}
\label{sec:felement_new}
Use transformation (\ref{tranformation:at2rz}), eq (\ref{eq:poisson7}) can be transformed into

\begin{equation}
\label{eq:poisson7-term1}
   \begin{array}{rl}
\mbox{term \#1}=
 & \ff{1}{\cal J} \pp{}{a} {\cal J} g_s \left( 
           \pp{\Phi}{a} g^{a a} + \pp{\Phi}{\theta} g^{a \theta}
                                        \right)
\\
 & \\
=& \ff{1}{\cal J} \pp{}{R} {\cal J} g_s \left( 
           \pp{\Phi}{a} g^{a a} + \pp{\Phi}{\theta} g^{a \theta}
                                        \right) \pp{R}{a}
   +
   \ff{1}{\cal J} \pp{}{Z} {\cal J} g_s \left( 
           \pp{\Phi}{a} g^{a a} + \pp{\Phi}{\theta} g^{a \theta}
                                        \right) \pp{Z}{a}
\\
 & \\
=&
   \ff{1}{\cal J} \pp{}{R} {\cal J} g_s \left[ 
           (\pp{\Phi}{R} \pp{R}{a} + \pp{\Phi}{Z} \pp{Z}{a}) g^{a a} + 
           (\pp{\Phi}{R} \pp{R}{\theta} +\pp{\Phi}{Z} \pp{Z}{\theta}) g^{a \theta}
                                        \right] \pp{R}{a} +
\\
 & \\
 & \ff{1}{\cal J} \pp{}{Z} {\cal J} g_s \left[ 
           (\pp{\Phi}{R} \pp{R}{a} + \pp{\Phi}{Z} \pp{Z}{a}) g^{a a} + 
           (\pp{\Phi}{R} \pp{R}{\theta} + \pp{\Phi}{Z} \pp{Z}{\theta}) g^{a \theta}
                                        \right] \pp{Z}{a}
\\
   \end{array}
\end{equation}

\begin{equation}
\label{eq:poisson7-term2}
   \begin{array}{rl}
\mbox{term \#2}= 
 & \ff{1}{\cal J} \pp{}{\theta} {\cal J} g_s \left(
           \pp{\Phi}{a} g^{a \theta} + \pp{\Phi}{\theta} g^{\theta \theta}
                                          \right)
\\
 & \\
=& \ff{1}{\cal J} \pp{}{R} {\cal J} g_s \left(
           \pp{\Phi}{a} g^{a \theta} + \pp{\Phi}{\theta} g^{\theta \theta}
                                          \right)\pp{R}{\theta} 
   +
   \ff{1}{\cal J} \pp{}{Z} {\cal J} g_s \left(
           \pp{\Phi}{a} g^{a \theta} + \pp{\Phi}{\theta} g^{\theta \theta}
                                          \right)\pp{Z}{\theta} 
\\
 & \\
=& \ff{1}{\cal J} \pp{}{R} {\cal J} g_s \left[
           (\pp{\Phi}{R} \pp{R}{a} + \pp{\Phi}{Z} \pp{Z}{a}) g^{a \theta} +
           (\pp{\Phi}{R} \pp{R}{\theta} + \pp{\Phi}{Z} \pp{Z}{\theta}) g^{\theta \theta})
                                          \right]\pp{R}{\theta} +
\\
 & \\
 & \ff{1}{\cal J} \pp{}{Z} {\cal J} g_s \left[
           (\pp{\Phi}{R} \pp{R}{a} + \pp{\Phi}{Z} \pp{Z}{a}) g^{a \theta} +
           (\pp{\Phi}{R} \pp{R}{\theta} + \pp{\Phi}{Z} \pp{Z}{\theta}) g^{\theta \theta})
                                          \right]\pp{Z}{\theta} 
\\
   \end{array}
\end{equation}

\begin{equation}
\label{eq:poisson7-term3}
   \begin{array}{rl}
\mbox{term \#3}= 
&           \ff{g_s}{q^2} \psp{\Phi}{\theta} g^{\varphi \varphi} =
            \ff{g_s}{q^2} g^{\varphi \varphi} \pp{}{\theta} \pp{\Phi}{\theta} \\
 & \\
=&
            \ff{g_s}{q^2} g^{\varphi \varphi} (\pp{}{R} \pp{\Phi}{\theta}) \pp{R}{\theta} +
            \ff{g_s}{q^2} g^{\varphi \varphi} (\pp{}{Z} \pp{\Phi}{\theta}) \pp{Z}{\theta}
\\
 & \\
=&
 \ff{g_s}{q^2} g^{\varphi \varphi} [\pp{}{R} (\pp{\Phi}{R} \pp{R}{\theta} + \pp{\Phi}{Z} \pp{Z}{\theta})] \pp{R}{\theta} +
 \ff{g_s}{q^2} g^{\varphi \varphi} [\pp{}{Z} (\pp{\Phi}{R} \pp{R}{\theta} + \pp{\Phi}{Z} \pp{Z}{\theta})] \pp{Z}{\theta}
   \end{array}
\end{equation}

\noindent
With finite elements $N_i$, they become

\begin{equation}
\label{eq:poisson7-term1}
   \begin{array}{rl}
\mbox{term \#1}=
& (\ff{1}{\cal J} \pp{R}{a} N_i) \pp{}{R} {\cal J} g_s \left[ 
           (\pp{N_j}{R} \pp{R}{a}      + \pp{N_j}{Z} \pp{Z}{a}) g^{a a} + 
           (\pp{N_j}{R} \pp{R}{\theta} + \pp{N_j}{Z} \pp{Z}{\theta}) g^{a \theta}
                                        \right] j +
\\
& \\
& (\ff{1}{\cal J} \pp{Z}{a} N_i) \pp{}{Z} {\cal J} g_s \left[ 
           (\pp{N_j}{R} \pp{R}{a}      + \pp{N_j}{Z} \pp{Z}{a}) g^{a a} + 
           (\pp{N_j}{R} \pp{R}{\theta} + \pp{N_j}{Z} \pp{Z}{\theta}) g^{a \theta}
                                        \right] j
\\
& \\
=&
 -\pp{}{R} (\ff{1}{\cal J} \pp{R}{a} N_i)  {\cal J} g_s \left[ 
           (\pp{N_j}{R} \pp{R}{a}      + \pp{N_j}{Z} \pp{Z}{a}) g^{a a} + 
           (\pp{N_j}{R} \pp{R}{\theta} + \pp{N_j}{Z} \pp{Z}{\theta}) g^{a \theta}
                                        \right] j +
\\
& \\
&-\pp{}{Z} (\ff{1}{\cal J} \pp{Z}{a} N_i)  {\cal J} g_s \left[ 
           (\pp{N_j}{R} \pp{R}{a}      + \pp{N_j}{Z} \pp{Z}{a}) g^{a a} + 
           (\pp{N_j}{R} \pp{R}{\theta} + \pp{N_j}{Z} \pp{Z}{\theta}) g^{a \theta}
                                        \right] j
\\
& \\
\sim&
 -\pp{N_i}{R}\pp{N_j}{R} (g_s \pp{R}{a} \pp{R}{a} g^{a a}) j
 -\pp{N_i}{R}\pp{N_j}{Z} (g_s \pp{R}{a} \pp{Z}{a} g^{a a}) j
 -\pp{N_i}{R}\pp{N_j}{R} (g_s \pp{R}{a} \pp{R}{\theta} g^{a \theta}) j
 -\pp{N_i}{R}\pp{N_j}{Z} (g_s \pp{R}{a} \pp{Z}{\theta} g^{a \theta}) j
\\
& \\
&
 -\pp{N_i}{Z}\pp{N_j}{R} (g_s \pp{Z}{a} \pp{R}{a} g^{a a}) j
 -\pp{N_i}{Z}\pp{N_j}{Z} (g_s \pp{Z}{a} \pp{Z}{a} g^{a a}) j
 -\pp{N_i}{Z}\pp{N_j}{R} (g_s \pp{Z}{a} \pp{R}{\theta} g^{a \theta}) j
 -\pp{N_i}{Z}\pp{N_j}{Z} (g_s \pp{Z}{a} \pp{Z}{\theta} g^{a \theta}) j
\\
& \\
=&
 -\pp{N_i}{R}\pp{N_j}{R} \left[g_s \pp{R}{a} (\pp{R}{a} g^{a a} + \pp{R}{\theta} g^{a \theta}) j \right]
 -\pp{N_i}{R}\pp{N_j}{Z} \left[g_s \pp{R}{a} (\pp{Z}{a} g^{a a} + \pp{Z}{\theta} g^{a \theta}) j \right]
\\
& \\
&
 -\pp{N_i}{Z}\pp{N_j}{R} \left[g_s \pp{Z}{a} (\pp{R}{a} g^{a a} + \pp{R}{\theta} g^{a \theta}) j \right]
 -\pp{N_i}{Z}\pp{N_j}{Z} \left[g_s \pp{Z}{a} (\pp{Z}{a} g^{a a} + \pp{Z}{\theta} g^{a \theta}) j \right]
\\
& \\
\equiv&
 -\pp{N_i}{R}\pp{N_j}{R} (\mbox{coef\_rr})
 -\pp{N_i}{R}\pp{N_j}{Z} (\mbox{coef\_rz})
 -\pp{N_i}{Z}\pp{N_j}{R} (\mbox{coef\_zr})
 -\pp{N_i}{Z}\pp{N_j}{Z} (\mbox{coef\_zz})
   \end{array}
\end{equation}

\begin{equation}
\label{eq:poisson7-term2}
   \begin{array}{rl}
\mbox{term \#2}= 
& (\ff{1}{\cal J} \pp{R}{\theta} N_i) \pp{}{R} {\cal J} g_s \left[
           (\pp{N_j}{R} \pp{R}{a}      + \pp{N_j}{Z} \pp{Z}{a}) g^{a \theta} +
           (\pp{N_j}{R} \pp{R}{\theta} + \pp{N_j}{Z} \pp{Z}{\theta}) g^{\theta \theta})
                                          \right]+
\\
& \\
& (\ff{1}{\cal J} \pp{Z}{\theta} N_i) \pp{}{Z} {\cal J} g_s \left[
           (\pp{N_j}{R} \pp{R}{a}      + \pp{N_j}{Z} \pp{Z}{a}) g^{a \theta} +
           (\pp{N_j}{R} \pp{R}{\theta} + \pp{N_j}{Z} \pp{Z}{\theta}) g^{\theta \theta})
                                          \right]
\\
& \\
= & 
 -\pp{}{R} (\ff{1}{\cal J} \pp{R}{\theta} N_i) {\cal J} g_s \left[
           (\pp{N_j}{R} \pp{R}{a}      + \pp{N_j}{Z} \pp{Z}{a}) g^{a \theta} +
           (\pp{N_j}{R} \pp{R}{\theta} + \pp{N_j}{Z} \pp{Z}{\theta}) g^{\theta \theta})
                                          \right]+
\\
& \\
&-\pp{}{Z} (\ff{1}{\cal J} \pp{Z}{\theta} N_i) {\cal J} g_s \left[
           (\pp{N_j}{R} \pp{R}{a}      + \pp{N_j}{Z} \pp{Z}{a}) g^{a \theta} +
           (\pp{N_j}{R} \pp{R}{\theta} + \pp{N_j}{Z} \pp{Z}{\theta}) g^{\theta \theta})
                                          \right]
\\
& \\
\sim & 
 -\pp{N_i}{R}\pp{N_j}{R} (g_s \pp{R}{\theta} \pp{R}{a} g^{a \theta})
 -\pp{N_i}{R}\pp{N_j}{Z} (g_s \pp{R}{\theta} \pp{Z}{a} g^{a \theta})
 -\pp{N_i}{R}\pp{N_j}{R} (g_s \pp{R}{\theta} \pp{R}{\theta} g^{\theta \theta})
 -\pp{N_i}{R}\pp{N_j}{Z} (g_s \pp{R}{\theta} \pp{Z}{\theta} g^{\theta \theta})
\\
& \\
&-\pp{N_i}{Z}\pp{N_j}{R} (g_s \pp{Z}{\theta} \pp{R}{a} g^{a \theta})
 -\pp{N_i}{Z}\pp{N_j}{Z} (g_s \pp{Z}{\theta} \pp{Z}{a} g^{a \theta})
 -\pp{N_i}{Z}\pp{N_j}{R} (g_s \pp{Z}{\theta} \pp{R}{\theta} g^{\theta \theta})
 -\pp{N_i}{Z}\pp{N_j}{Z} (g_s \pp{Z}{\theta} \pp{Z}{\theta} g^{\theta \theta})
\\
& \\
=&
 -\pp{N_i}{R}\pp{N_j}{R} \left[g_s \pp{R}{\theta} (\pp{R}{a} g^{a \theta} + \pp{R}{\theta} g^{\theta \theta}) \right]
 -\pp{N_i}{R}\pp{N_j}{Z} \left[g_s \pp{R}{\theta} (\pp{Z}{a} g^{a \theta} + \pp{Z}{\theta} g^{\theta \theta}) \right]
\\
& \\
&-\pp{N_i}{Z}\pp{N_j}{R} \left[g_s \pp{Z}{\theta} (\pp{R}{a}  g^{a \theta} + \pp{R}{\theta} g^{\theta \theta}) \right]
 -\pp{N_i}{Z}\pp{N_j}{Z} \left[g_s \pp{Z}{\theta} (\pp{Z}{a}) g^{a \theta} + \pp{Z}{\theta} g^{\theta \theta}) \right]
\\
& \\
\equiv &
 -\pp{N_i}{R}\pp{N_j}{R} \mbox{coef\_rr}
 -\pp{N_i}{R}\pp{N_j}{Z} \mbox{coef\_rz}
 -\pp{N_i}{Z}\pp{N_j}{R} \mbox{coef\_zr}
 -\pp{N_i}{Z}\pp{N_j}{Z} \mbox{coef\_zz}
   \end{array}
\end{equation}

\begin{equation}
\label{eq:poisson7-term3}
   \begin{array}{rl}
\mbox{term \#3}= 
&
 \ff{g_s}{q^2} g^{\varphi \varphi} [\pp{}{R} (\pp{\Phi}{R} \pp{R}{\theta} + \pp{\Phi}{Z} \pp{Z}{\theta})] \pp{R}{\theta} +
 \ff{g_s}{q^2} g^{\varphi \varphi} [\pp{}{Z} (\pp{\Phi}{R} \pp{R}{\theta} + \pp{\Phi}{Z} \pp{Z}{\theta})] \pp{Z}{\theta}
\\
=& (\ff{g_s}{q^2} \pp{R}{\theta} g^{\varphi \varphi} N_i) [\pp{}{R} (\pp{N_j}{R} \pp{R}{\theta} + \pp{N_j}{Z} \pp{Z}{\theta})] +
   (\ff{g_s}{q^2} \pp{Z}{\theta} g^{\varphi \varphi} N_i) [\pp{}{Z} (\pp{N_j}{R} \pp{R}{\theta} + \pp{N_j}{Z} \pp{Z}{\theta})]
\\
=&-\pp{}{R} (\ff{g_s}{q^2} \pp{R}{\theta} g^{\varphi \varphi} N_i) (\pp{N_j}{R} \pp{R}{\theta} + \pp{N_j}{Z} \pp{Z}{\theta})
  -\pp{}{Z} (\ff{g_s}{q^2} \pp{Z}{\theta} g^{\varphi \varphi} N_i) (\pp{N_j}{R} \pp{R}{\theta} + \pp{N_j}{Z} \pp{Z}{\theta})
\\
\sim & -\pp{N_i}{R} (\ff{g_s}{q^2} \pp{R}{\theta} g^{\varphi \varphi} ) (\pp{N_j}{R} \pp{R}{\theta} + \pp{N_j}{Z} \pp{Z}{\theta})
       -\pp{N_i}{Z} (\ff{g_s}{q^2} \pp{Z}{\theta} g^{\varphi \varphi} ) (\pp{N_j}{R} \pp{R}{\theta} + \pp{N_j}{Z} \pp{Z}{\theta})
\\
=& -\pp{N_i}{R}\pp{N_j}{R} (\ff{g_s}{q^2} g^{\varphi \varphi} \pp{R}{\theta} \pp{R}{\theta})
   -\pp{N_i}{R}\pp{N_j}{Z} (\ff{g_s}{q^2} g^{\varphi \varphi} \pp{R}{\theta} \pp{Z}{\theta})
   -\pp{N_i}{Z}\pp{N_j}{R} (\ff{g_s}{q^2} g^{\varphi \varphi} \pp{Z}{\theta} \pp{R}{\theta})
   -\pp{N_i}{Z}\pp{N_j}{Z} (\ff{g_s}{q^2} g^{\varphi \varphi} \pp{Z}{\theta} \pp{Z}{\theta})
\\
=& -\pp{N_i}{R}\pp{N_j}{R} \mbox{coef\_rr}
   -\pp{N_i}{R}\pp{N_j}{Z} \mbox{coef\_rz}
   -\pp{N_i}{Z}\pp{N_j}{R} \mbox{coef\_zr}
   -\pp{N_i}{Z}\pp{N_j}{Z} \mbox{coef\_zz}
   \end{array}
\end{equation}


% ................................... %
\subsection{Boundary Conditon}
\label{sec:bc}
The solution domain is bounded by two magnetic flux surfaces.
One is the innermost flux surface $a=a0$ and the other one is the outermost flux surface $a=a1$.
Dirichlet boundary condition and Neumann boundary condition can be specified on
either one of these two surfaces.
We use a set of array to store the index of the grid points which sit on boudnaries.

The Dirichlet boundary condition can be imposed by
zeroing the row of the matrix with indices on the boundaries and then set their diagonal 
entries to $1$.
The rhs vector will be changed acoordingly by setting those elements to the real solution.
Here is a simple example.
Assumping there are $5$ grid points and they are numbered as $1, 2, 3, 4, 5$.
The linear system is given as
$$
\left[ \begin{array}{lllll}
         a1 & a2 & a3 & a4 & a5 \\
         b1 & b2 & b3 & b4 & b5 \\
         c1 & c2 & c3 & c4 & c5 \\
         d1 & d2 & d3 & d4 & d5 \\
         e1 & e2 & e3 & e4 & e5
\end{array} \right]
\left[ \begin{array}{l}
         x1\\
         x2\\
         x3\\
         x4\\
         x5
\end{array} \right]
=
\left[ \begin{array}{l}
         b1\\
         b2\\
         b3\\
         b4\\
         b5
\end{array} \right]
$$
Suppose the boundary points are at $2$ and $5$, and the Dirichlet bounnday condition is given on
these points, i.e.
$$
\begin{array}{l}
x2=v1\\
x5=v2\\
\end{array}
$$
Then the matrix and rhs are modified in the following way
$$
\left[ \begin{array}{lllll}
         a1 & a2 & a3 & a4 & a5 \\
            &  1 &    &    &    \\
         c1 & c2 & c3 & c4 & c5 \\
         d1 & d2 & d3 & d4 & d5 \\
            &    &    &    &  1
\end{array} \right]
\left[ \begin{array}{l}
         x1\\
         x2\\
         x3\\
         x4\\
         x5
\end{array} \right]
=
\left[ \begin{array}{l}
         b1\\
         v2\\
         b3\\
         v4\\
         v5
\end{array} \right]
$$
By solving the above linear system, we immediately get $x2=v1$ and $x5=v2$.
This is unsymmetrical zeroing.
We can do symmetric zeroing to keep the modified matrix symmetrical
so that better linear solver can be applied. It can be explained 
by continuing zero the corresponding columns, i.e., column $2$ and $5$,
of the row-zeroing modified matrix
$$
\left[ \begin{array}{lllll}
         a1 & a2-a2 & a3 & a4 & a5-a5 \\
            &  1    &    &    &       \\
         c1 & c2-c2 & c3 & c4 & c5-c5 \\
         d1 & d2-d2 & d3 & d4 & d5-d5 \\
            &       &    &    &  1
\end{array} \right]
\left[ \begin{array}{l}
         x1\\
         x2\\
         x3\\
         x4\\
         x5
\end{array} \right]
=
\left[ \begin{array}{l}
         b1-a2v1-a5v2\\
         v2\\
         b3-c2v1-c5v2\\
         v4-d2v1-d5v2\\
         v5
\end{array} \right]
$$
i.e.,
$$
\left[ \begin{array}{lllll}
         a1 &    & a3 & a4 &   \\
            &  1 &    &    &   \\
         c1 &    & c3 & c4 &   \\
         d1 &    & d3 & d4 &   \\
            &    &    &    &  1
\end{array} \right]
\left[ \begin{array}{l}
         x1\\
         x2\\
         x3\\
         x4\\
         x5
\end{array} \right]
=
\left[ \begin{array}{l}
         b1-a2v1-a5v2\\
         v2\\
         b3-c2v1-c5v2\\
         v4-d2v1-d5v2\\
         v5
\end{array} \right]
$$
Note that the rhs vector $b$ was modified acoordingly.

\noindent
If Neumann boundary condition is imposed, the correction is much more complicated.
To make the explanation easier, we write the problem as follows
$$
\begin{array}{l}
\Delta u = f \\
\pp{u}{n} = g \\
\end{array}
$$
and the numerical linear system as
$$A x = b $$
We start from the variational form 
$$
\int\int N_i \bm{\nabla}_\perp \cdot g_s \bm{\nabla}_\perp \Phi {\cal J} dad\theta
=
\int {\cal J} g_s N_i \bm{n} \cdot \bm{\nabla}_\perp \Phi dl
- \int\int \bm{\nabla}_\perp({\cal J} N_i) \cdot g_s \bm{\nabla}_\perp \Phi dad\theta
$$
Suppose on the boundary we have
$$
\bm{n} \cdot \bm{\nabla}_\perp \Phi = \gamma(a,\theta)
$$
$ \bm{n} $ is the boundary normal directon pointing away from the solution domain.
In the current version of GTS, this directon is equivalent to $\ff{\bm{\nabla}a}{|\bm{\nabla}a|}$.
Thus,
$$
\ff{\bm{\nabla}a}{|\bm{\nabla}a|} \cdot \bm{\nabla} \Phi = \gamma(a,\theta)
$$
Using the previous expression for $\bm{\nabla} \Phi$, it becomes
$$
\ff{\bm{\nabla}a}{|\bm{\nabla}a|} \cdot 
(\pp{\Phi}{a} \bm{\nabla}a + \pp{\Phi}{\theta} \bm{\nabla}\theta) 
= \gamma(a,\theta)
$$
i.e.,
$$
\ff{1}{|\bm{\nabla}a|} \left( g^{aa} \pp{\Phi}{a} + g^{a\theta}\pp{\Phi}{\theta} \right) = \gamma(a,\theta)
$$
Therefore
$$
\begin{array}{l}
\int {\cal J} g_s N_i \bm{n} \cdot \bm{\nabla}_\perp dl \\
= \int {\cal J} g_sN_i \ff{1}{|\bm{\nabla}a|} ( g^{aa} \pp{\Phi}{a} + g^{a\theta}\pp{\Phi}{\theta} ) dl \\
= \int_{e_1} {\cal J} g_s N_i \ff{1}{|\bm{\nabla}a|} ( g^{aa} \pp{\Phi}{a} + g^{a\theta}\pp{\Phi}{\theta} ) dl
  +
  \int_{e_2} {\cal J} g_s N_i \ff{1}{|\bm{\nabla}a|} ( g^{aa} \pp{\Phi}{a} + g^{a\theta}\pp{\Phi}{\theta} ) dl
  +
  \int_{e_3} {\cal J} g_s N_i \ff{1}{|\bm{\nabla}a|} ( g^{aa} \pp{\Phi}{a} + g^{a\theta}\pp{\Phi}{\theta} ) dl
\\
 \equiv I_1 + I_2 + I3 \\
= \int {\cal J} g_s N_i \ff{1}{|\bm{\nabla}a|} \gamma(a,\theta) dl
\end{array}
$$
where $I_1$, $I_2$, and $I_3$ are the linear integral along one of the 3 edges of each triangle,
respectively.
For simplicity, we use the $C^0$ linear elements $N_i$.

(1) if the boundary falls on edge $1$,
$$
\begin{array}{l}
N_1=1-\xi\\
N_2=\xi\\
N_3=0 \\
a = \sum_i a_i N_i = a_1(1-\xi) + a_2 \xi = a_1 + \bar{a}_{21} \xi\\
\theta = \sum_i \theta_i N_i = \theta_1(1-\xi) + \theta_2 \xi = \theta_1 + \bar{\theta}_{21} \xi\\
dl^2 = da^2 + d\theta^2 
     = \bar{a}_{21}^2 d\xi^2 + \bar{\theta}_{21}^2 d\xi^2
     = (\bar{a}_{21}^2 + \bar{\theta}_{21}^2) d\xi^2
     = l_{21}^2 d\xi^2
\end{array}
$$
The following terms are used to modify the right hand side
$$
\begin{array}{ll}
rhs(v_1) & = rhs(v_1) - \int_0^1 N_1(\xi,0) {\cal J} g_s \ff{1}{|\bm{\nabla}a|} \gamma l_{21} d\xi \\
         & = rhs(v_1) - \int_0^1 (1-\xi) {\cal J} g_s \ff{1}{|\bm{\nabla}a|} \gamma l_{21} d\xi \\
         & \simeq rhs(v_1) - \ff{1}{2} {\cal J} g_s \ff{1}{|\bm{\nabla}a|} \gamma l_{21}
\\
rhs(v_2) & = rhs(v_2) - \int_0^1 N_2(\xi,0) {\cal J} g_s \ff{1}{|\bm{\nabla}a|} \gamma l_{21} d\xi \\
         & = rhs(v_2) - \int_0^1 (\xi) {\cal J} g_s \ff{1}{|\bm{\nabla}a|} \gamma l_{21} d\xi \\
         & \simeq rhs(v_2) - \ff{1}{2} {\cal J} g_s \ff{1}{|\bm{\nabla}a|} \gamma l_{21}
\\
rhs(v_3) & = rhs(v_3) - \int_0^1 N_3(\xi,0) {\cal J} g_s \ff{1}{|\bm{\nabla}a|} \gamma l_{21} d\xi \\
         & = rhs(v_3) - \int_0^1 (0) {\cal J} g_s \ff{1}{|\bm{\nabla}a|} \gamma l_{21} d\xi \\
         & = rhs(v_3)
\end{array}
$$

(2) if the boundary falls on edge $2$,
$$
\begin{array}{l}
N_1=1-\xi-\eta=0\\
N_1=1-\eta\\
N_2=\eta\\
a = \sum_i a_i N_i = a_2(1-\eta) + a_3 \eta = a_2 + \bar{a}_{32} \xi\\
\theta = \sum_i \theta_i N_i = \theta_2(1-\eta) + \theta_3 \eta = \theta_2 + \bar{\theta}_{32} \xi\\
dl^2 = da^2 + d\theta^2 
     = \bar{a}_{32}^2 d\eta^2 + \bar{\theta}_{32}^2 d\eta^2
     = (\bar{a}_{32}^2 + \bar{\theta}_{32}^2) d\eta^2
     = l_{32}^2 d\eta^2
\end{array}
$$
The following terms are used to modify the right hand side
$$
\begin{array}{ll}
rhs(v_1) & = rhs(v_1) - \int_0^1 N_1(\xi,\eta) {\cal J} g_s \ff{1}{|\bm{\nabla}a|} \gamma l_{32} d\eta \\
         & = rhs(v_1) - \int_0^1 (0) {\cal J} g_s \ff{1}{|\bm{\nabla}a|} \gamma l_{32} d\eta \\
         & = rhs(v_1)
\\
rhs(v_2) & = rhs(v_2) - \int_0^1 N_2(\xi,\eta) {\cal J} g_s \ff{1}{|\bm{\nabla}a|} \gamma l_{32} d\eta \\
         & = rhs(v_2) - \int_0^1 (1-\eta) {\cal J} g_s \ff{1}{|\bm{\nabla}a|} \gamma l_{32} d\eta \\
         & \simeq rhs(v_2) - \ff{1}{2} {\cal J} g_s \ff{1}{|\bm{\nabla}a|} \gamma l_{32}
\\
rhs(v_3) & = rhs(v_3) - \int_0^1 N_3(\xi,\eta) {\cal J} g_s \ff{1}{|\bm{\nabla}a|} \gamma l_{32} d\eta \\
         & = rhs(v_3) - \int_0^1 (\eta) {\cal J} g_s \ff{1}{|\bm{\nabla}a|} \gamma l_{32} d\eta \\
         & \simeq rhs(v_3) - \ff{1}{2} {\cal J} g_s \ff{1}{|\bm{\nabla}a|} \gamma l_{32}
\end{array}
$$

(3) if the boundary falls on edge $3$,
$$
\begin{array}{l}
N_1=1-\eta\\
N_2=0\\
N_3=\eta
a = \sum_i a_i N_i = a_1(1-\eta) + a_3 \eta = a_1 + \bar{a}_{13} \eta\\
\theta = \sum_i \theta_i N_i = \theta_1(1-\eta) + \theta_3 \eta = \theta_1 + \bar{\theta}_{31} \eta\\
dl^2 = da^2 + d\theta^2 
     = \bar{a}_{13}^2 d\xi^2 + \bar{\theta}_{13}^2 d\eta^2
     = (\bar{a}_{13}^2 + \bar{\theta}_{13}^2) d\eta^2
     = l_{13}^2 d\eta^2
\end{array}
$$
The following terms are used to modify the right hand side
$$
\begin{array}{ll}
rhs(v_1) & = rhs(v_1) - \int_0^1 N_1(0,\eta) {\cal J} g_s \ff{1}{|\bm{\nabla}a|} \gamma (-l_{13}) d\eta \\
         & = rhs(v_1) - \int_0^1 (1-\eta) {\cal J} g_s \ff{1}{|\bm{\nabla}a|} \gamma (-l_{13}) d\eta \\
         & \simeq rhs(v_1) - \ff{1}{2} {\cal J} g_s \ff{1}{|\bm{\nabla}a|} \gamma (-l_{13})
\\
rhs(v_2) & = rhs(v_2) - \int_0^1 N_2(0,\eta) {\cal J} g_s \ff{1}{|\bm{\nabla}a|} \gamma (-l_{13}) d\eta \\
         & = rhs(v_2) - \int_0^1 (0) {\cal J} g_s \ff{1}{|\bm{\nabla}a|} \gamma (-l_{13}) d\eta \\
         & = rhs(v_2)
\\
rhs(v_3) & = rhs(v_3) - \int_0^1 N_3(0,\eta) {\cal J} g_s \ff{1}{|\bm{\nabla}a|} \gamma (-l_{13}) d\eta \\
         & = rhs(v_3) - \int_0^1 (\eta) {\cal J} g_s \ff{1}{|\bm{\nabla}a|} \gamma (-l_{13}) d\eta \\
         & \simeq rhs(v_3) - \ff{1}{2} {\cal J} g_s \ff{1}{|\bm{\nabla}a|} \gamma (-l_{13})
\end{array}
$$
These are implemented in subroutine $neumbcfix$.


The stiffness matrix $\bm{A}$ is singular when neumann bc is specified.
Thus the nullspace of this matrix is not empty,
and in almost every case the constant vector $\bm{e}=\{1\}$ belongs to this space,
i.e.,

\begin{verbatim}
   nullspace=span{e}
\end{verbatim}

\noindent
To check if this is true, we just need to confirm the following the matrix-vector product of $A$ and $\bm{e}$

\begin{verbatim}
   A.e = 0  
\end{verbatim}

\noindent
is zero. Then We need to check the solvability by computing the vector product of 
rhs vector $\bm{b}$ and the constant vector $\bm{e}$

\begin{verbatim}
   b.e
\end{verbatim}

\noindent
If dot product is zero, then there is a solution;
if not, the following projection is needed

\begin{verbatim}
  b = b - e (b.e)/(e.e)
\end{verbatim}

\noindent
to ensure that there exists a solution in the linear system we are solving.
If a solution exists, it’s unique though only up to an additive constant;
in this case there’s also a requirement on $\gamma$.
This follows from $\int\int \Delta u d\sigma = \int\int f d\sigma$. With partial integration,
$\int \pp{u}{n} dl = \int\int f d\sigma$, i.e., 
$\int \gamma dl = \int\int f d\sigma$.

\noindent
Here is the implementation after the above steps

\begin{verbatim}
  /* declare nullspace and create nullspace object*/
  MatNullSpace  nullsp;
  PetscBool  isNull;
       call MatNullSpaceCreate(partd_comm, PETSC_TRUE, 0, PETSC_NULL_OBJECT, nullsp, ierr);
       call MatNullSpaceTest(nullsp, Amat, isNull,ierr)
       call KSPSetNullSpace(ksp, nullsp,ierr)
       call KSPGetPC(ksp,prec,ierr)
       call PCFactorSetShiftType(prec,MAT_SHIFT_POSITIVE_DEFINITE,ierr)
       call PCFactorSetShiftAmount(prec,1.e-10,ierr)
       call MatNullSpaceDestroy(nullsp,ierr)
\end{verbatim}

\noindent
After it is solved, the solution vector $\bm{x}$ needs to be projected away from the nullspace by
\begin{verbatim}
  x = x - e (x.e)/(e.e)
\end{verbatim}

\noindent
in the end.

When matrix is not big, we can use the following technique. First an
$n$-dimension space can be spanned by
$$
R^n=span{e,q2,q3,\cdots,qn}
$$
where $q2, q3, \cdots, qn$ are $n-1$ linear independent vectors and 
they forming the following transformation matrix
$$
Q=\left[ e,q2,q3,\cdots,qn \right]
$$
Apply this transformation to our linear system $Ax=b$, we get
$$(Q^T A Q) (Q^T x) = Q^T b
$$
$$
(Q^T A Q) = \left[ \begin{array}{ll}  
                         0 & 0 \\
                         0 & B \\
                   \end{array}  
            \right],
(Q^T B) = \left[ \begin{array}{l}  
                         0 \\
                         d \\
                   \end{array}  
            \right],
(Q^T x) = \left[ \begin{array}{l}  
                         0 \\
                         z \\
                   \end{array}  
            \right]
$$
So instead of solving $Ax=b$, we solve $Bz=d$. $x$ is recovered from
$$
x = Q \left[ \begin{array}{l}  
                         w \\
                         z \\
                   \end{array}  
            \right]
$$
where $w$ is arbitrary constant.
If we take $w=0$, we get the least square solution.

The example of dealing with Neumann bc is given in subroutine $testsolve$
by setting $whichone=4$ and $bc\_type=0$.
$$
u=a^2, \pp{u}{a}=2a, \pp{u}{t}=0
$$

$$
\pp{u}{n}
=\bm{n} \cdot \bm{\nabla} u
=\ff{\bm{\nabla}a}{|\bm{\nabla}a|} \cdot \bm{\nabla} u
=\ff{\bm{\nabla}a}{|\bm{\nabla}a|} \cdot \pp{u}{a} \bm{\nabla}a
=\ff{1}{|\bm{\nabla}a|} \pp{u}{a} g^{aa}
\equiv \gamma(a)
$$



% ................................... %
\subsection{Numerical Difficulty: periodic domain}
\label{sec:period}
The numerical difficulty comes from the $2\pi$ peroidicity in $\theta$ direction
when numerical integration is performed.
The value of $2\pi$ must be subtracted from $\theta$ 
whenever the elements has one or two nodes crossing the $\theta=0, 2\pi$ line.
This is implemented in subroutine $fixtheta$.



% ................................... %
\begin{thebibliography}{1}

  \bibitem{WW1} W. X. Wang et al, {\em Gyro-kinetic silumation of global turbulent
                transport properties in tokamak experiments},
                Physcs of  Plasma, vol 13,  2006

  \bibitem{Coor} W. D. D'haeseleer, {\em Flux coordinates and magnetic field structure},
                1991.

  \bibitem{Integration} G. E. Karniadakis and S. J. Sherwin, {\em Spectral/hp Element Methods for CFD},
                1999.

  \bibitem{C1element1} Gouri Dhatt, Emmanuel Lefrançois, Gilbert Touzot, {\em Finite Element Method}, 2012

  \bibitem{C1element0} A. R. Mitchell and R. Wait, Wiley, {\em The Finite Element Method in Partial Differential Equations}, London, 1977
  \end{thebibliography}



% ................................... %
\begin{appendices}
\section{Flux Coordiante Basics}
\label{App:coor}
We review some important facts here about curved coordiante system 
that will be needed in the the derivation of pde equations and 
in the implementation of finite elements in the flux coordiantes.
We list them here so we will not mention them when we use these facts.

As a start, we introduce the coordinate surface $u^i=\mbox{const}$.
$\bm{\nabla} u^i$ is perpendicular to the $i$th surface.
$i=1,2,3$.

The coordinate curve $u^i$
is defined by the intersection of $2$ surfaces 
$u^j=\mbox{const}$, $u^k=\mbox{const}$.
On this curve, $du^j = du^k = 0$, $j\neq k, j\neq i, k\neq i$.

$$
\bm{R} = \bm{R} (u^1, u^2, u^3),
\;\;\;
\bm{e}^i =\bm{\nabla}u^i,
\;\;\;
\bm{e}_i = \pp{\bm{R}}{u^i}
\;\;\;
\bm{e}^i \cdot \bm{e}_i = \delta_j^i
$$

$$
\bm{e}^i = \ff{ \bm{e}_j \times \bm{e}_k }{ \bm{e}_i \cdot \bm{e}_j \times \bm{e}_k }
         = \ff{1}{\cal J} \bm{e}_j \times \bm{e}_k
$$

$$
\bm{e}_i = \ff{\bm{R}}{u^i}
         = \ff{ \bm{e}^j \times \bm{e}^k }{ \bm{e}^i \cdot \bm{e}^j \times \bm{e}^k }
         = {\cal J} \bm{e}^j \times \bm{e}^k
$$

$$
D_i = \bm{D} \cdot \bm{e}_i,
D^j = \bm{D} \cdot \bm{e}^j,
\bm{D} = \sum_i D_i \bm{e}^i = \sum_j D^j \bm{e}_j 
$$

$du^i$ are contravariants.
$\bm{\nabla}\Phi$ is a covariant vector: 
$$\bm{\nabla}\Phi = \pp{\Phi}{u^i} \bm{\nabla}u^i = \pp{\Phi}{u^i} \bm{e}^i$$

1). gradient
$$\bm{\nabla}\Phi = \bm{\nabla}u^i \pp{\Phi}{u^i} 
                  = \pp{\Phi}{u^i} \bm{e}^i
$$

2). divergence
$$
\bm{\nabla} \cdot \bm{D} = \ff{1}{\cal J} \pp{}{u^i} \left( {\cal J} D^i \right)
                         = \ff{1}{\cal J} \pp{}{u^i} \left( {\cal J} \bm{D} \cdot \bm{e}^i \right)
                         = \ff{1}{\cal J} \pp{}{u^i} \left( {\cal J} D_j \bm{e}^j \cdot \bm{e}^i \right)
                         = \ff{1}{\cal J} \pp{}{u^i} \left( {\cal J} D_j g^{ji} \right)
$$

3). laplacian
$$
\Delta D = \ff{1}{\cal J} \pp{}{u^i} \left( {\cal J} \bm{\nabla} D \cdot \bm{e}^i \right)
         = \ff{1}{\cal J} \pp{}{u^i} \left( {\cal J} \pp{D}{u^j} \bm{e}^j \cdot \bm{e}^i \right)
         = \ff{1}{\cal J} \pp{}{u^i} \left( {\cal J} \pp{D}{u^j} g^{ji} \right)
$$

4). Jacobian
$$
{\cal J} = \pp{\Phi}{u^1} \cdot \pp{\Phi}{u^2} \times \pp{\Phi}{u^3}
$$

5). metrics
$$
g^{aa} = \bm{\nabla} a \cdot \bm{\nabla} a,
\;\;\;
g^{a\theta} = \bm{\nabla} a \cdot \bm{\nabla} \theta,
\;\;\;
g^{\theta\theta} = \bm{\nabla} \theta \cdot \bm{\nabla} \theta,
$$

6). the differential arc length
$$
dl(\mbox{along} \; u^j) 
\equiv dl(i)
= \| d\bm{R} \|
= \sqrt{ d\bm{R}(i) \cdot d\bm{R}(i)}
$$
while 
$$d\bm{R}(i) = \pp{\bm{R}}{u^1}du^1
               + \pp{\bm{R}}{u^2}du^2
               + \pp{\bm{R}}{u^3}du^3
             = \bm{e}_i du^i
             = {\cal J} \| \bm{\nabla}u^j \times \bm{\nabla}u^j \| du^i
$$

7). the differential area element
$$
\begin{array}{ll}
dS (\mbox{in } u^i
 \equiv \mbox{const})
& \equiv dS(i) \\
& = \| d\bm{R}(j) \times \bm{R}(k) \| \\
& = \| \pp{\bm{R}}{u^j} \times \pp{\bm{R}}{u^k} \| du^j du^k \\
& = \| \bm{e}_j \times \bm{e}_k \| du^j du^k \\
& = {\cal J} \| \bm{e}^i \| du^j du^k \\
& = {\cal J} \| \bm{\nabla} u^i \| du^j du^k
\end{array}
$$



% ................................... %
\section{Numerical integration}
\label{App:gauss}
The basic problem considered by numerical integration is to compute a definite integral
\begin{equation}
\label{eq:int1}
\int u(\xi) d\xi
\end{equation}
There are many methods of approximating the integral with arbitrary precision.
A method 
which yields a small error for a small number of evaluations 
is usually considered superior. 
Gauss Quadrature is one of such algorithms.

The integrand is evaluated at a finite set of points called integration points 
and a weighted sum of these values is used to approximate the integral. 
The integration points and weights depend on the specific method used 
and the accuracy required from the approximation.

\iffalse
\begin{verbatim}
~/project/spectralElement/QUADRATURE/stroud.f90 
~/project/spectralElement/GLpts/quadrature_lib.f 
therm2d used "rule=1"
lagrange used "rule=2"
10532 !
10533 !  1 point, precision 1.
10534 !
10535   if ( rule == 1 ) then
10536 
10537     a = 1.0D+00 / 3.0D+00
10538     w = 1.0D+00
10539 
10540     norder = 1
10541     xtab(1) = a
10542     ytab(1) = a
10543     weight(1) = w
10544 !
10545 !  3 points, precision 2, Strang and Fix formula #1.
10546 !
10547   else if ( rule == 2 ) then
10548 
10549     a = 1.0D+00
10550     b = 3.0D+00
10551     c = 4.0D+00
10552     d = 6.0D+00
10553 
10554     norder = 3
10555     xtab(1:3) =   (/ c, a, a /) / d
10556     ytab(1:3) =   (/ a, c, a /) / d
10557     weight(1:3) = (/ a, a, a /) / b
\end{verbatim}
\fi

In Gauss quadrature technique, the integrand is represented as a Lagrange polynomial
using $Q$ points $\xi_i$
$$
u(\xi) = \sum_{i=0}^{Q-1} u(\xi_i) h_i(\xi) + \epsilon(u)
$$
If we substitute it into the integral (\ref{eq:int1}), we obtain
a representation of the integral as a summation
$$
\int u(\xi) d\xi = \sum_{i=0}^{Q-1} w_i u(\xi_i) + R(u)
$$
where
$$
\begin{array}{l}
w_i = \int_{-1}^1 h_i(\xi) d\xi \\
R(u) = \int_{-1}^1 \epsilon(u) d\xi
\end{array}
$$
Since $u(\xi)$ is represented by a polynomial of order $Q-1$,
we would expect the above relation to be exact if $u(\xi)$ 
is a polynomial of order $Q-1$ or less.
This would be true if the points are equispaced in the interval.
There is, however, a better choice of abscissae which permits exact 
integration of polynomials of higher order than $Q-1$.

If we ensure that the quadrature formula is exact for polynomials of order $k$
at least, then the formula is also exact for the next $(k-1)$ order of polynomials.




\end{appendices}
\end{document}
\iffalse
\fi
\begin{verbatim}
\end{verbatim}
